%% (H1) Společná matematika
\subsection{1. Základy diferenciálního a integrálního počtu}
\begin{itemize}
\item Posloupnosti reálných čísel a jejich limity (definice)
\begin{itemize}
\item posloupnost s hodnotami v množině $M$ je zobrazení z $\mathbb N$ do $M$
\item nechť $(a_n)$ je reálná posloupnost a $L \in \mathbb R^*$
\item pokud $(\forall\varepsilon\gt 0)(\exists n_0\in\mathbb N):(n\geq n_0\implies |a_n-L|\lt\varepsilon)$, píšeme, že $\lim a_n=L$ nebo $\lim_{n\to\infty}a_n=L$ nebo $a_n\to L$, a řekneme, že posloupnost $(a_n)$ má limitu $L$
\item pro $L\in \mathbb R$ mluvíme o vlastní limitě a pro $L=\pm\infty$ o limitě nevlastní
\item posloupnost, která má vlastní limitu, konverguje (pokud ji nemá, tak diverguje)
\item každá posloupnost reálných čísel má nejvýše jednu limitu
\end{itemize}
\item Aritmetika limit
\begin{itemize}
\item nechť $(a_n),(b_n)$ jsou posloupnosti reálných čísel takové, že $\lim a_n=a$ a $\lim b_n=b$
\item pak\dots{}
\begin{itemize}
\item $\lim\,(a_n+b_n)=a+b$
\item $\lim\,(a_nb_n)=ab$
\item pokud $b_n\neq 0$ pro každé $n\gt 0$, pak $\lim\,(a_n/b_n)=a/b$
\end{itemize}
\item to vše platí za předpokladu, že je výraz na pravé straně definován
\begin{itemize}
\item neurčité výrazy (nejsou definované): součet nekonečen s různými znaménky, součin nekonečna s nulou, podíl nekonečen, dělení reálného čísla nulou
\end{itemize}
\item pokud je $(a_n)$ omezená a $(b_n)$ konverguje k nule, pak $\lim\,(a_nb_n)=0$
\begin{itemize}
\item tzn. limita $(a_n)$ nemusí existovat
\end{itemize}
\end{itemize}
\item Limity a uspořádání
\begin{itemize}
\item nechť $(a_n),(b_n)$ jsou posloupnosti reálných čísel takové, že mají limity $\lim a_n=a$ a $\lim b_n=b$
\item když $a\lt b$, tak existuje $n_0$, že $n\gt n_0\implies a_n\lt b_n$
\item a naopak, když $a_n\leq b_n$ pro každé $n\gt n_0$, pak $a\leq b$
\item poznámka: může se stát, že $\forall n:a_n\lt b_n$, ale $\lim a_n=\lim b_n$
\begin{itemize}
\item např. uvažujme $a_n=1-\frac1n\lt b_n=1$, přičemž $\lim a_n=\lim b_n=1$
\end{itemize}
\end{itemize}
\item Věta o dvou policajtech
\begin{itemize}
\item nechť posloupnosti $(a_n),(b_n),(c_n)\subseteq\mathbb R$ splňují
\begin{itemize}
\item $\lim a_n=\lim b_n=a\in\mathbb R$
\item $\forall n\gt n_0:a_n\leq c_n\leq b_n$
\end{itemize}
\item pak $(c_n)$ konverguje a $\lim c_n=a$
\item platí to i pro nevlastní limity - tam nám stačí jeden policajt
\end{itemize}
\item Součet řady a částečný součet řady
\begin{itemize}
\item nekonečná řada (reálných čísel) je výraz \[\sum_{n=1}^\infty a_n=a_1+a_2+a_3+\dots\]
\begin{itemize}
\item kde $(a_n)_{n\geq 1}$ je posloupnost reálných čísel
\end{itemize}
\item ($n$-tý) částečný součet řady se značí $s_n$ a je to součet jejích prvních $n$ členů ($s_n=a_1+a_2+\dots+a_n$)
\item pokud existuje vlastní limita posloupnosti $(s_n)$ částečných součtů dané řady, mluvíme o \textit{konvergentní} řadě a limita $\lim s_n$ je jejím \textit{součtem}
\item pokud $\lim s_n$ neexistuje nebo je nevlastní, je daná řada \textit{divergentní}
\end{itemize}
\item Geometrická řada, harmonická řada
\begin{itemize}
\item geometrická řada
\begin{itemize}
\item $\sum_{n=0}^\infty q^n=1+q+q^2+q^3+\dots$
\item $q\in\mathbb R$ je parametr zvaný \textit{kvocient}
\item pro součet geometrické řady platí \[\sum_{n=0}^\infty q^n=\begin{cases}\frac1{1-q} &-1\lt q\lt1\\ +\infty &q\geq 1\\\text{neexistuje}&q\leq -1\end{cases}\]
\item vyplývá to ze vzorce pro částečný součet $s_n=\frac{q^n-1}{q-1}$
\end{itemize}
\item harmonická řada
\begin{itemize}
\item $\sum_{n=1}^\infty\frac 1n=1+\frac12+\frac13+\dots$
\item diverguje, přestože její prvky konvergují k nule
\end{itemize}
\item pro řady tvaru $\sum_{n=1}^\infty\frac1{n^s}$ obecně platí, že konvergují pro $s\gt 1$ a divergují pro $s\leq 1$
\end{itemize}
\item Limita funkce v bodě: definice, aritmetika limit
\begin{itemize}
\item okolí bodu je interval $U(a,\delta)=(a-\delta,a+\delta)$
\begin{itemize}
\item okolí nekonečen definujeme jako $U(+\infty,\delta)=(1/\delta,+\infty)$, podobně pro $-\infty$
\item definujeme také pravé okolí bodu $U^+(a,\delta)=[a,a+\delta)$, podobně levé okolí $U^-$
\item prstencová okolí bodu $a\in\mathbb R$ definujeme jako obyčejná okolí s vyjmutým bodem $a$, značí se jako $P(a,\delta)$ apod.
\end{itemize}
\item řekneme, že funkce $f$ má v bodě $a\in\mathbb R^*$ limitu $A\in\mathbb R$, když $(\forall\varepsilon\gt 0)(\exists\delta\gt 0):x\in P(a,\delta)\implies f(x)\in U(A,\varepsilon)$
\begin{itemize}
\item zapisujeme $\lim_{x\to a}f(x)=A$
\item poznámka: limita funkce $f$ v bodě $a$ nezávisí na její hodnota v $a$, funkce $f$ dokonce ani nemusí být v $a$ definovaná
\item pokud místo prstencového okolí $P$ použijeme pravé prstencové okolí $P^+$, dostaneme definici limity zprava, značí se $\lim_{x\to a^+}f(x)$
\begin{itemize}
\item obdobně lze definovat limitu zleva
\item pokud $\lim_{x\to a^+}f(x)=\lim_{x\to a^-}f(x)=A$, pak $\lim_{x\to a}f(x)=A$
\item naopak pokud jsou limity zprava a zleva různé, limita neexistuje
\end{itemize}
\item podobně jako u posloupnosti - pokud limita funkce existuje, je jednoznačně určena
\item limitu funkce lze rovněž definovat pomocí posloupnosti
\end{itemize}
\item aritmetika limit
\begin{itemize}
\item nechť $a,A,B\in\mathbb R^*$, nechť $f,g$ jsou funkce definované na nějakém prstencovém okolí $P(a,\delta)$ bodu $a$, nechť platí $\lim_{x\to a}f(x)=A$ a $\lim_{x\to a}g(x)=B$
\item pak platí $\lim_{x\to a}f(x)+g(x)=A+B$, je-li tento součet definován
\item $\lim_{x\to a}f(x)g(x)=AB$, je-li tento součin definován
\item pokud je $g(x)$ nenulová na nějakém prstencovém okolí bodu $a$, pak $\lim_{x\to a}f(x)/g(x)=A/B$, je-li tento podíl definován
\end{itemize}
\end{itemize}
\item Limita funkce v bodě: vztah s uspořádáním
\begin{itemize}
\item nechť $c\in\mathbb R^*$ a funkce $f,g,h$ jsou definované na prstencovém okolí bodu $c$
\item mají-li $f,g$ v bodě $c$ limitu a $\lim_{x\to c}f(x)\gt\lim_{x\to c}g(x)$, pak existuje $\delta\gt 0$ takové, že $f(x)\gt g(x)$ pro každé $x\in P(c,\delta)$
\item existuje-li $\delta\gt 0$ takové, že $f(x)\geq g(x)$ pro každé $x\in P(c,\delta)$, a mají-li funkce $f,g$ limitu v bodě $c$, potom $\lim_{x\to c}f(x)\geq \lim_{x\to c}g(x)$
\item existuje-li $\delta\gt 0$ takové, že $f(x)\leq h(x)\leq g(x)$ pro každé $x\in P(c,\delta)$ a $\lim_{x\to c}f(x)=\lim_{x\to c}g(x)=A\in\mathbb R^*$ potom i $\lim_{x\to c}h(x)=A$
\begin{itemize}
\item v podstatě dva policajti
\end{itemize}
\end{itemize}
\item Limita funkce v bodě: limita složené funkce
\begin{itemize}
\item nechť $A,B,C\in\mathbb R^*$, nechť $g(x)$ je funkce splňující $\lim_{x\to A}g(x)=B$ a nechť $f(x)$ je funkce splňující $\lim_{x\to B}f(x)=C$
\item navíc nechť je splněna aspoň jedna z následujících podmínek
\begin{itemize}
\item $f(x)$ je spojitá v $B$
\item na nějakém prstencovém okolí bodu $A$ funkce $g(x)$ nikde nenabývá hodnotu $B$
\end{itemize}
\item potom $\lim_{x\to A}f(g(x))=C$
\end{itemize}
\item Spojitost a spojitost na intervalu
\begin{itemize}
\item řekneme, že funkce $f$ je v bodě $a\in\mathbb R$ spojitá, pokud $\lim_{x\to a}f(x)=f(a)$
\begin{itemize}
\item spojitost zprava/zleva definujeme pomocí jednostranných limit
\end{itemize}
\item funkce je spojitá na intervalu, pokud je spojitá v každém vnitřním bodu a pokud je (odpovídajícím způsobem) jednostranně spojitá v každém krajním bodu
\item funkce může být v daném bodě buď spojitá, nebo nespojitá, nebo nedefinovaná
\end{itemize}
\item Funkce spojité na intervalu: nabývání mezihodnot
\begin{itemize}
\item nechť $a\lt b$ jsou reálná čísla
\item nechť $f:[a,b]\to\mathbb R$ je na intervalu $[a,b]$ spojitá
\item označme $m=\min\set{f(a),f(b)}$ a $M=\max\set{f(a),f(b)}$
\item pak každé reálné číslo z intervalu $[m,M]$ je hodnotou funkce $f$ (pro každé $y\in[m,M]$ existuje $x\in[a,b]$ takové, že $f(x)=y$)
\end{itemize}
\item Funkce spojité na intervalu: nabývání maxima
\begin{itemize}
\item nechť $a\leq b$ jsou reálná čísla
\item nechť $f:[a,b]\to\mathbb R$ je spojitá funkce
\item potom $f$ nabývá na intervalu $[a,b]$ svého maxima (i minima)
\end{itemize}
\item Derivace - definice a základní pravidla pro výpočet
\begin{itemize}
\item nechť $f:M\to\mathbb R$, $b\in M$ a $U(b,\delta)\subseteq M$ pro nějaké $\delta\gt 0$
\item derivace funkce $f$ v bodě $b$ je limita \[f'(b):=\lim_{h\to 0}\frac{f(b+h)-f(b)}{h}=\lim_{x\to b}\frac{f(x)-f(b)}{x-b}\]
\item také můžeme mít derivaci zprava/zleva, pak je limita jednostranná
\item derivace buď existuje vlastní, nebo existuje nevlastní, nebo neexistuje
\item pokud má $f$ v $b$ vlastní derivaci, říkáme, že $f$ je v $b$ diferencovatelná
\item pro pravidla derivací viz libovolnou vhodnou tabulku, některá složitější jsou v přípravě na zápočtový test
\begin{practicebox}
Poznámky zde derivační pravidla nerozepisují. Natrénuj derivace součinu, podílu a složené funkce na konkrétních příkladech (viz tabulka derivací).
\end{practicebox}
\end{itemize}
\item L'Hospitalovo pravidlo
\begin{itemize}
\item nechť $a\in\mathbb R^*$, nechť funkce $f,g:P(a,\delta)\to\mathbb R$ mají na $P(a,\delta)$ vlastní derivaci a nechť $g'(x)\neq 0$ na $P(a,\delta)$
\item pokud $\lim_{x\to a}f(x)=\lim_{x\to a}g(x)=0$ a $\lim_{x\to a}\frac{f'(x)}{g'(x)}=A\in\mathbb R^*$, pak i $\lim_{x\to a}\frac{f(x)}{g(x)}=A$
\item pokud $\lim_{x\to a} |g(x)|=+\infty$ a $\lim_{x\to a}\frac{f'(x)}{g'(x)}=A\in\mathbb R^*$, pak i $\lim_{x\to a}\frac{f(x)}{g(x)}=A$
\end{itemize}
\item Vyšetření průběhu funkcí: extrémy, monotonie a konvexita/konkavita
\begin{itemize}
\item extrémy
\begin{itemize}
\item pokud má funkce $f$ v bodě $a$ nenulovou derivaci $f'(a)\neq 0$, pak $f$ v $a$ nenabývá lokální extrém
\item tedy pokud má $f$ lokální extrém bodě $a$, tak buď $f'(a)=0$, nebo $f'(a)$ neexistuje
\item pozor, funkce může mít globální extrém v krajním bodě uzavřeného intervalu
\end{itemize}
\item monotonie
\begin{itemize}
\item uvažujeme interval s kladnou délkou a spojitou funkci $f$ s derivací v každém vnitřním bodě intervalu
\item pokud na vnitřku intervalu platí $f'\gt 0$, je $f$ na daném intervalu rostoucí
\item pokud $f'\geq 0$, pak je neklesající
\item pokud $f'\lt 0$, pak je klesající
\item pokud $f'\leq 0$, pak je nerostoucí
\item pokud $f'=0$, pak je konstantní
\end{itemize}
\item konvexita/konkavita
\begin{itemize}
\item zjednodušená definice: funkce $f$ je na intervalu $(a,b)$ konvexní, pokud graf funkce na intervalu $(a,b)$ leží pod úsečkou spojující body $(a,f(a))$ a $(b,f(b))$
\begin{itemize}
\item respektive pokud leží pod nebo na úsečce, tak je konvexní, pokud leží ostře pod úsečkou, tak je ryze konvexní
\item podobně konkávnost
\end{itemize}
\item pokud je druhá derivace funkce kladná, je ryze konvexní (pokud je nezáporná, je konvexní)
\item pokud je druhá derivace funkce záporná, je ryze konkávní (pokud je nekladná, je konkávní)
\end{itemize}
\end{itemize}
\item Taylorův polynom (limitní forma)
\begin{itemize}
\item Taylorův polynom řádu $n$ funkce $f$ v bodě $a$: \[T_n^{f,a}(x)=\sum_{i=0}^n\frac{f^{(i)}(a)}{i!}(x-a)^i\]
\begin{itemize}
\item poznámka: nultý člen součtu bude vždy $f(a)$
\end{itemize}
\item má-li funkce $f$ v bodě $a\in\mathbb R$ derivace všech řádů, rozumíme pro $x\in\mathbb R$ její taylorovou řadou se středem v $a$ řadu \[T^{f,a}(x)=\sum_{n=0}^\infty \frac{f^{(n)}(a)}{n!}(x-a)^n\]
\begin{itemize}
\item součet Taylorovy řady je $f(x)$
\end{itemize}
\end{itemize}
\item Primitivní funkce: definice a metody výpočtu (substituce, per-partes)
\begin{itemize}
\item nechť $-\infty\leq a\lt b\leq +\infty$ a $f:(a,b)\to\mathbb R$ je daná funkce
\item pokud má funkce $F:(a,b)\to\mathbb R$ na $(a,b)$ derivaci a ta se rovná $f(x)$, řekneme, že $F$ je na intervalu $(a,b)$ primitivní funkcí k funkci $f$
\item primitivní funkce je jednoznačná až na konstantu
\item Newtonův integrál funkce $f$ na intervalu $(a,b)$ definujeme jako \[(N)\int_a^bf(x)\,dx:=[F]_a^b=F(b^-)-F(a^+)=\lim_{x\to b^-}F(x)-\lim_{x\to a^+}F(x)\]
\begin{itemize}
\item poznámka: konstanta $c$ se odečte
\end{itemize}
\item substituci a per-partes je nutné nastudovat (i pro určitý integrál)
\begin{practicebox}
Výpočetní techniky (substituce, per partes) si projdi na příkladech - v poznámkách je jen vzorec, ne postup.
\end{practicebox}
\item vzorec pro per-partes: $\int u'v=uv-\int uv'$
\end{itemize}
\item Riemannův integrál: definice, souvislost s primitivní funkcí (Newtonovým integrálem)
\begin{itemize}
\item uvažujeme dělení $D=(a_0,a_1,\dots,a_k)$ intervalu $[a,b]$
\begin{itemize}
\item tzn. $a=a_0\lt a_1\lt a_2\lt \dots\lt a_k=b$
\end{itemize}
\item definujeme dolní Riemannovu sumu $s(f,D)=\sum_{i=0}^{k-1} |I_i|m_i$
\begin{itemize}
\item $I_i=[a_i,a_{i+1}]$
\item $m_i$ je infimum funkce $f$ v intervalu $I_i$
\end{itemize}
\item také definujeme horní Riemannovu sumu $S(f,D)=\sum_{i=0}^{k-1}|I_i|M_i$
\begin{itemize}
\item $M_i$ \dots{} supremum
\end{itemize}
\item dolní Riemannův integrál na intervalu $[a,b]$ definujeme jako supremum z $s(f,D)$ pro všechna dělení $D$ intervalu $[a,b]$
\begin{itemize}
\item značí se $\underline{\int_a^b f}$
\end{itemize}
\item horní Riemannův integrál je infimum z $S(f,D)$
\begin{itemize}
\item $\overline{\int_a^bf}$
\end{itemize}
\item funkce má Riemannův integrál, pokud se horní a dolní R. integrály rovnají (pak tuhle společnou hodnotu nazveme Riemannovým integrálem)
\begin{itemize}
\item $\int_a^b f$
\end{itemize}
\item 2\textbackslash{}. základní věta analýzy: pokud Riemannův a Newtonův integrál existují, pak jsou si rovny
\end{itemize}
\item Aplikace integrálů: odhady součtu řad (konečných i nekonečných)
\begin{itemize}
\item pro přirozené $n$ a neklesající funkci $f$ na intervalu $[1,n]$ platí \[\sum_{k=1}^{n-1}f(k)\leq\int_1^n f\leq\sum_{k=2}^n f(k)\]
\item integrální kritérium konvergence
\begin{itemize}
\item uvažujme nerostoucí nezápornou funkci $f$
\item řada $\sum_{k=1}^\infty f(k)$ konverguje $\iff\lim_{b\to+\infty}\int_1^b f\lt +\infty$
\end{itemize}
\item přesný postup pro odhad součtu řady je potřeba nacvičit
\begin{practicebox}
Odhad součtu řady přes integrál je výpočetní dovednost - vyřeš pár konkrétních příkladů.
\end{practicebox}
\end{itemize}
\item Aplikace integrálů: obsahy rovinných útvarů
\begin{itemize}
\item plocha rovinného útvaru $U(a,b,f)$ pod grafem funkce $f$ \dots{} $\int_a^b f$
\end{itemize}
\item Aplikace integrálů: objemy a povrchy rotačních útvarů v prostoru
\begin{itemize}
\item $V$ \dots{} rotační těleso vzniklé rotací rovinného útvaru $U(a,b,f)$ pod grafem funkce $f$ kolem osy $x$
\item $\mathrm{objem}(V)=\pi\int_a^b f(t)^2\,dt$
\item $\mathrm{povrch}(V)=2\pi\int_a^b f(t)\sqrt{1+(f'(t))^2}\,dt$
\end{itemize}
\item Aplikace integrálů: délka křivky
\begin{itemize}
\item uvažujeme křivku z bodu $a$ do bodu $b$ popsanou funkcí $f$
\item $\int_a^b\sqrt{1+(f'(t))^2}\,dt$
\end{itemize}
\end{itemize}
\subsection{2. Algebra a lineární algebra}
\begin{itemize}
\item Grupy a podgrupy, permutace
\begin{itemize}
\item binární operace na množině $X$ je zobrazení $X \times X \to X$
\item neutrální prvek: $(\exists e \in G)(\forall a \in G): a\circ e = e\circ a = a$
\item inverzní prvek: $(\forall a \in G)(\exists b \in G): a\circ b = b\circ a = e$
\item grupa $(G,\circ)$ je množina $G$ spolu s binární operací $\circ$ na $G$ splňující asociativitu operace $\circ$, existenci neutrálního prvku a existenci inverzních prvků
\item pokud je navíc operace $\circ$ komutativní, pak se jedná o abelovskou grupu
\item grupa $(H,\bullet)$ je podgrupou grupy $(G,\circ)$, pokud $H\subseteq G$ a $\forall a,b\in H:a\bullet b=a\circ b$
\begin{itemize}
\item píšeme $(H,\bullet)\leq (G,\circ)$
\end{itemize}
\item permutace na množině $\lbrace1,2,\dots,n\rbrace$ je bijektivní zobrazení $p:\lbrace1,2,\dots,n\rbrace\rightarrow \lbrace1,2,\dots,n\rbrace$
\begin{itemize}
\item skládání permutací \dots{} $(q\circ p)(i)=q(p(i))$
\item permutační matice
\begin{itemize}
\item $(P)_{i,j} = \begin{cases} 1 &\text{pokud } p(i)=j \\ 0 &\text{jinak}\end{cases}$
\end{itemize}
\item cyklus \dots{} např. $(1,2,3,4)$
\begin{itemize}
\item odpovídá zobrazení, které mapuje 1$\to$2, 2$\to$3, 3$\to$4, 4$\to$1
\end{itemize}
\item transpozice \dots{} permutace, která má pouze jeden netriviální cyklus o délce 2, tedy např. $(1,3)$
\item jakoukoliv permutaci lze rozložit na transpozice
\begin{itemize}
\item cyklus $(1,2,3,4)$ lze rozložit na $(1,4)\circ(1,3)\circ(1,2)$ nebo na $(1,2)\circ(2,3)\circ(3,4)$
\end{itemize}
\item znaménko permutace $p$ je $\text{sgn}(p)=(-1)^{\text{po\text{č}et inverz\text{í} v}\ p}$
\begin{itemize}
\item inverze v $p$ je dvojice prvků $(i,j):i \lt j \land p(i) \gt p(j)$
\item nebo také $\text{sgn}(p)=(-1)^{\text{po\text{č}et transpozic v}\ p}$
\end{itemize}
\end{itemize}
\end{itemize}
\item Tělesa a speciálně konečná tělesa
\begin{itemize}
\item těleso je množina $\mathbb K$ spolu se dvěma komutativními binárními operacemi $+$ a $\cdot$, kde $(\mathbb K, +)$ a $(\mathbb K \setminus \lbrace0\rbrace, \cdot)$ jsou abelovské grupy a navíc platí distributivita $\forall a,b,c \in \mathbb K : a\cdot (b+c)=(a\cdot b)+(a\cdot c)$
\item charakteristika tělesa
\begin{itemize}
\item v tělese $\mathbb K$, pokud $\exists n \in \mathbb N: \underbrace{1+1+\dots+1}_{n}=0$, pak nejmenší takové $n$ je charakteristika tělesa $\mathbb K$
\item jinak má těleso $\mathbb K$ charakteristiku 0
\item značí se $\text{char}(\mathbb K)$
\item lze dokázat, že u konečných těles musí být prvočíselná
\end{itemize}
\item $\mathbb Z_p$ je těleso, právě když je $p$ prvočíslo
\item ale existují i konečná tělesa s neprvočíselným řádem (počtem prvků) - např. Galois field $GF(4)$
\begin{itemize}
\item řád musí být kladná celá mocnina prvočísla
\end{itemize}
\end{itemize}
\item Soustavy lineárních rovnic - maticový zápis, elementární řádkové úpravy, odstupňovaný tvar matice
\begin{itemize}
\item maticový zápis soustavy rovnic
\begin{itemize}
\item soustava $m$ lineárních rovnic o $n$ neznámých \dots{} $Ax = b$
\item rozšířená matice soustavy \dots{} $(A|b) \in \mathbb{R}^{m\times (n+1)}$
\item matice soustavy \dots{} $A \in \mathbb{R}^{m\times n}$
\item vektor pravých stran \dots{} $b \in \mathbb{R}^m$
\item vektor neznámých \dots{} $x = (x_1,\dots,x_n)^T$
\item vektor $x \in \mathbb{R}^n$ je řešení soustavy $Ax = b$, pokud splňuje všechny její rovnice
\item soustavy $Ax = 0$ se nazývají homogenní a vždy umožňují $x = 0$
\end{itemize}
\item elementární řádkové operace
\begin{itemize}
\item definujeme základní dvě elementární řádkové úpravy
\begin{itemize}
\item vynásobení $i$-tého řádku nenulovým $t \in \mathbb R \setminus \lbrace0\rbrace$
\item přičtení $j$-tého řádku k $i$-tému řádku
\end{itemize}
\item z těch lze odvodit další dvě úpravy
\begin{itemize}
\item přičtení $t$-násobku $j$-tého řádku k $i$-tému řádku ($t$ může být i nulové)
\item záměna dvou řádků
\end{itemize}
\item provedení jedné elementární úpravy značíme $A\sim A'$
\item provedení posloupnosti úprav značíme $A\sim\sim A'$
\item elementárními úpravami soustavy $(A|b)$ se nemění množina řešení
\end{itemize}
\item odstupňovaný tvar matice
\begin{itemize}
\item matice je v řádkově odstupňovaném tvaru (REF = row echelon form), pokud jsou nenulové řádky (ostře) seřazeny podle počtu počátečních nul a nulové řádky jsou pod nenulovými
\begin{itemize}
\item k formální definici bychom zavedli notaci $j(i)$ pro pozici prvního nenulového prvku v $i$-tém řádku
\end{itemize}
\item první nenulový prvek nenulového řádku se nazývá pivot, pod pivotem jsou v REF všechny prvky nulové
\item pro soustavu $A'x = b'$ s $A'$ v REF jsou proměnné odpovídající sloupcům s pivoty bázické, ostatní jsou volné
\end{itemize}
\item redukovaný odstupňovaný tvar (RREF)
\begin{itemize}
\item pivoty jsou jednotkové
\item nad i pod pivoty jsou nuly
\end{itemize}
\end{itemize}
\item Gaussova a Gaussova-Jordanova eliminace, popis množiny řešení
\begin{itemize}
\item Gaussova eliminace = převod do odstupňovaného tvaru pomocí elementárních úprav, to se dá popsat algoritmicky:
\begin{itemize}
\item seřadíme řádky podle počtu počátečních nul
\item najdeme první dva řádky, co mají stejně počátečních nul
\item od toho druhého (a všech dalších, co mají stejně nul) odečteme vhodný násobek toho prvního, abychom mu první nenulu vynulovali
\item tohle opakujeme, dokud se nedostaneme do REF
\end{itemize}
\item Gaussova-Jordanova eliminace = převod do RREF
\begin{itemize}
\item z pivotů se jedničky dělají snadno - stačí řádek vydělit pivotem
\item nad pivoty se jedničky také dostávají snadno - stačí od řádku odečíst vhodný násobek řádku s pivotem v daném sloupci (dává smysl postupovat odspodu)
\end{itemize}
\item množina řešení
\begin{itemize}
\item pokud je pivot v posledním sloupci rozšířené matice soustavy, soustava rovnic nemá řešení
\item existuje bijekce $\bar x \mapsto \bar x + x^0$ mezi množinami $\lbrace \bar x: A\bar x=0\rbrace$ a $\lbrace x: Ax=b\rbrace$
\begin{itemize}
\item pokud $x^0$ splňuje $Ax^0=b$
\end{itemize}
\item množinu řešení soustavy $Ax=b$ popíšeme jako $x=y+p_1\overline x_1+p_2\overline x_2+\dots+p_{n-r}\overline x_{n-r}$
\item $y$ \dots{} libovolné řešení soustavy $Ax=b$
\begin{itemize}
\item u homogenní soustavy to typicky bude nulový vektor
\item je to člen, který odpovídá tomu $x^0$ výše
\end{itemize}
\item $p_i\overline x_i$
\begin{itemize}
\item takový člen tam bude za každou volnou proměnnou (bude jich $n-r$, tedy počet řádků - hodnost matice)
\item $p_i$ je libovolný reálný parametr
\item tyhle členy můžeme získat tak, že řešíme soustavu $A\overline x=0$ a řešení vyjádříme pomocí parametrů (a pak je vytkneme)
\item nebo je lze vyčíst z RREF matice
\item ve vektoru $\overline x_i$ bude typicky jednička na místě odpovídající volné proměnné a nuly u všech ostatních volných proměnných (protože je to vlastně složka, která se stará o konkrétní volnou proměnnou)
\end{itemize}
\end{itemize}
\end{itemize}
\item Operace s maticemi a základní typy matic, hodnost matice
\begin{itemize}
\item hodnost matice
\begin{itemize}
\item hodnost matice $A$, značená jako $\mathrm{rank}(A)$, je počet pivotů v libovolné $A'$ v REF takové, že $A\sim\sim A'$
\end{itemize}
\item jednotková matice
\begin{itemize}
\item pro $n \in \mathbb{N}$ je jednotková matice $I_n \in \mathbb{R}^{n\times n}$ definovaná tak, že $(I_n)_{i,j} = 1 \iff i=j$, ostatní prvky jsou nulové 
\end{itemize}
\item transponovaná matice
\begin{itemize}
\item transponovaná matice k matici $A \in \mathbb{R}^{m\times n}$ je matice $A^T \in \mathbb R^{n\times m}$ splňující $(A^T)_{i,j}=a_{j,i}$
\end{itemize}
\item symetrická matice
\begin{itemize}
\item čtvercová matice A je symetrická, pokud $A^T =A$, tedy $a_{i,j}=a_{j,i}$
\end{itemize}
\item maticový součin
\begin{itemize}
\item pro $A \in \mathbb R^{m\times n},B\in \mathbb R^{n\times p}$ je součin $(AB) \in \mathbb R^{m\times p}$ definován $(AB)_{i,j}=\sum^{n}_{k=1}a_{i,k}b_{k,j}$
\end{itemize}
\end{itemize}
\item Regulární a inverzní matice
\begin{itemize}
\item inverzní matice
\begin{itemize}
\item pokud pro čtvercovou matici $A \in \mathbb R^{n\times n}$ existuje $B \in \mathbb R^{n\times n}$ taková, že $AB=I_n$, pak se $B$ nazývá inverzní matice a značí se $A^{-1}$
\item výpočet: $(A|I_n)\sim\sim (I_n|A^{-1})$
\end{itemize}
\item regulární/singulární matice
\begin{itemize}
\item pokud má matice $A$ inverzi, pak se nazývá regulární, jinak je singulární
\end{itemize}
\item věta: pro čtvercovou matici $A \in \mathbb R^{n\times n}$ jsou následující podmínky ekvivalentní
\begin{enumerate}
\item matice A je regulární, tedy k ní existuje inverzní matice \dots{} $\exists B: AB=I_n$
\item $\text{rank}(A) = n$
\item $A\sim\sim I_n$
\item systém $Ax = 0$ má pouze triviální řešení $x = 0$
\end{enumerate}
\item inverzní matici k matici $A$ lze najít pomocí Gaussovy-Jordanovy eliminace $(A|I_n)\sim\sim(I_n|A^{-1})$
\begin{itemize}
\item pokud tento proces selže, tak je $A$ singulární
\end{itemize}
\end{itemize}
\item Vektorový prostor, lineární kombinace, lineární závislost a nezávislost, lineární obal, systém generátorů
\begin{itemize}
\item vektorový prostor
\begin{itemize}
\item vektorový prostor $(V,+,\cdot)$ nad tělesem $(\mathbb K, +,\cdot)$ je množina spolu s binární operací $+$ na $V$ a binární operací skalárního násobku $\cdot: \mathbb K \times V \rightarrow V$
\item $(V,+)$ je abelovská grupa
\item $\forall \alpha,\beta \in \mathbb K, \forall u,v \in V$
\begin{itemize}
\item asociativita \dots{} $(\alpha \cdot \beta) \cdot u = \alpha \cdot (\beta \cdot u)$
\item neutrální prvek (skalár) vůči násobení skalárem \dots{} $1 \cdot u = u$
\item distributivita \dots{} $(\alpha + \beta) \cdot u = (\alpha \cdot u) + (\beta \cdot u)$
\item distributivita \dots{} $\alpha \cdot (u+v)=(\alpha \cdot u)+(\alpha \cdot v)$
\end{itemize}
\item prvky $\mathbb K$ se nazývají skaláry, prvky $V$ vektory
\item rozlišujeme nulový skalár $0$ a nulový vektor $o$
\end{itemize}
\item podprostor vektorového prostoru
\begin{itemize}
\item nechť $V$ je vektorový prostor na $\mathbb K$, potom podprostor $U$ je neprázdná podmnožina $V$ splňující uzavřenost na součet vektorů a uzavřenost na násobení skalárem (z $\mathbb K$) - z toho nutně vyplývá $o \in U$
\end{itemize}
\item lineární kombinace
\begin{itemize}
\item lineární kombinace vektorů $v_1,\dots,v_k \in V$ nad $\mathbb K$ je libovolný vektor $u = \alpha_1v_1+\dots+\alpha_kv_k$, kde $\alpha_1,\dots,\alpha_k \in \mathbb K$
\end{itemize}
\item lineárně nezávislé vektory
\begin{itemize}
\item množina vektorů $X$ je lineárně nezávislá, pokud nulový vektor nelze získat netriviální lineární kombinací vektorů z $X$; v ostatních případech je množina $X$ lineárně závislá
\item vektory $v_1,\dots,v_n$ jsou lineárně nezávislé $\equiv$ $\sum_{i=1}^n \alpha_iv_i=o \iff \alpha_1=\dots=\alpha_n=0$
\end{itemize}
\item lineární obal (podprostor generovaný množinou)
\begin{itemize}
\item lineární obal $\mathcal L(X)$ podmnožiny $X$ vektorového prostoru $V$ je průnik všech podprostorů $U$ z $V$, které obsahují $X$
\item alternativní značení: $\mathrm{span}(X)$
\item pro $X \subseteq V$ platí $\text{span}(X)=\bigcap U:U\Subset V, X\subseteq U$
\item jde o podprostor generovaný $X$, vektory v množině $X$ se označují jako generátory podprostoru
\end{itemize}
\item věta: průnik podprostorů je taky podprostor
\item věta: $\mathcal L(X)$ je množina všech lineárních kombinací vektorů z $X$
\end{itemize}
\item Steinitzova věta o výměně, báze, dimenze, souřadnice
\begin{itemize}
\item lemma o výměně
\begin{itemize}
\item buď $y_1,\dots,y_n$ systém generátorů vektorového prostoru $V$
\item nechť vektor $x \in V$ má vyjádření $x = \sum_{i=1}^n\alpha_iy_i$
\item pak pro libovolné $k$ takové, že $\alpha_k \neq 0$, je $y_1,\dots,y_{k-1},x,y_{k+1},\dots,y_n$ systém generátorů prostoru $V$
\end{itemize}
\item Steinitzova věta o výměně
\begin{itemize}
\item buď $V$ vektorový prostor
\item buď $x_1,\dots,x_m$ lineárně nezávislý systém ve $V$
\item nechť $y_1,\dots,y_n$ je systém generátorů $V$
\item pak platí $m\leq n$ a existují navzájem různé indexy $k_1,\dots,k_{n-m}$ takové, že $x_1,\dots,x_m,y_{k_1},\dots,y_{k_{n-m}}$ tvoří systém generátorů $V$
\end{itemize}
\item báze vektorového prostoru
\begin{itemize}
\item báze vektorového prostoru $V$ je lineárně nezávislá množina $X$, která generuje $V$ (tedy $\text{span}(X)=V$)
\end{itemize}
\item dimenze vektorového prostoru
\begin{itemize}
\item dimenze konečně generovaného vektorového prostoru $V$ je mohutnost (kardinalita / počet prvků) kterékoli z jeho bází; značí se $\mathrm{dim}(V)$
\end{itemize}
\item vektor souřadnic
\begin{itemize}
\item nechť $X=(v_1,\dots,v_n)$ je uspořádaná báze vektorového prostoru $V$ nad $\mathbb K$, potom vektor souřadnic $u \in V$ vzhledem k bázi $X$ je $[u]_x=(\alpha_1,\dots,\alpha_n)^T \in \mathbb K^n$, kde $u=\sum_{i=1}^n\alpha_iv_i$
\end{itemize}
\end{itemize}
\item Vektorové podprostory, zejména maticové (řádkový, sloupcový, jádro) a jejich dimenze
\begin{itemize}
\item řádkový a sloupcový prostor matice $A \in \mathbb K^{m\times n}$
\begin{itemize}
\item sloupcový prostor $\mathcal S(A) \subseteq \mathbb K^m$ je lineární obal sloupců $A$
\item řádkový prostor $\mathcal R(A) \subseteq \mathbb K^n$ je lineární obal řádků $A$
\item $\mathcal S(A)=\lbrace u \in \mathbb K^m:u=Ax,\,x\in \mathbb K^n \rbrace$
\item $\mathcal R(A)=\lbrace v \in \mathbb K^n:v=A^Ty,\,y\in \mathbb K^m \rbrace$
\end{itemize}
\item jádro matice $A \in \mathbb K^{m\times n}$
\begin{itemize}
\item $\text{ker}(A) = \lbrace x \in \mathbb K^n: Ax=0\rbrace$
\end{itemize}
\item věta: $\forall A \in \mathbb K^{m\times n}:\text{dim}(\text{ker}(A))+\text{rank}(A)=n$
\item pozorování: dimenze řádkového a sloupcového prostoru se shodují (a rovnají se hodnosti)
\end{itemize}
\item Lineární zobrazení - definice, maticová reprezentace lineárního zobrazení, matice složeného zobrazení
\begin{itemize}
\item lineární zobrazení
\begin{itemize}
\item nechť $U$ a $V$ jsou vektorové prostory nad stejným tělesem $\mathbb K$
\item zobrazení $f:U\rightarrow V$ nazveme lineární, pokud splňuje $\forall u,v \in U, \forall \alpha \in \mathbb K:$
\begin{itemize}
\item $f(u+v)=f(u)+f(v)$
\item $f(\alpha\cdot u)=\alpha \cdot f(u)$
\begin{itemize}
\item z toho vyplývá, že pro lineární zobrazení obecně platí $f(o) = o$
\end{itemize}
\end{itemize}
\item věta: lineární zobrazení je jednoznačně definováno tím, kam zobrazí vektory báze (proto můžeme definovat matici lineárního zobrazení)
\end{itemize}
\item matice lineárního zobrazení
\begin{itemize}
\item nechť $U$ a $V$ jsou vektorové prostory nad stejným tělesem $\mathbb K$ s bázemi $X=(u_1,\dots,u_n)$ a $Y=(v_1,\dots,v_m)$
\item matice lineárního zobrazení $f:U\rightarrow V$ vzhledem k bázím $X$ a $Y$ je $[f]_{X,Y} \in \mathbb K^{m\times n}$, jejíž sloupce jsou vektory souřadnic obrazů vektorů báze $X$ vzhledem k bázi $Y$, tedy $[f(u_1)]_Y,\dots,[f(u_n)]_Y$
\item pro $w \in U$ tedy platí, že $[f(w)]_Y=[f]_{X,Y}[w]_X$
\end{itemize}
\item matice složeného zobrazení
\begin{itemize}
\item mějme vektorové prostory $U,V,W$ s konečnými uspořádanými bázemi $B,C,D$
\item pro matice lineárních zobrazení $f:U\to V$ a $g:V\to W$ platí vztah $[g\circ f]_{B,D}=[g]_{C,D}[f]_{B,C}$
\item je to hezčí, když použijeme Hladíkovo značení: ${}_D[g\circ f]_B={}_D[g]_C\cdot{}_C[f]_B$
\end{itemize}
\item matice přechodu
\begin{itemize}
\item nechť $X$ a $Y$ jsou dvě konečné báze vektorového prostoru $U$
\item matice přechodu od $X$ k $Y$ je $[id]_{X,Y}$
\item pro $u \in U$ tedy platí, že $[u]_Y=[id(u)]_Y=[id]_{X,Y}[u]_X$
\item matice přechodu je regulární, platí $[id]_{Y,X}=([id]_{X,Y})^{-1}$
\item výpočet: $[id]_{X,Y}=Y^{-1}X$ nebo také $(Y|X)\sim\sim(I_n|[id]_{X,Y})$
\end{itemize}
\end{itemize}
\item Obraz a jádro lineárních zobrazení
\begin{itemize}
\item nechť $U$ a $V$ jsou vektorové prostory nad stejným tělesem $\mathbb K$
\item jádro lineárního zobrazení
\begin{itemize}
\item $\text{ker}(f)=\lbrace w \in U:f(w)=o\rbrace$
\item vektory v prostoru $U$, které se zobrazí na nulový vektor v prostoru $V$
\end{itemize}
\item obraz $f(U)$ podprostoru $U$ je podprostorem $V$
\begin{itemize}
\item pokud máme matici zobrazení, tak obraz získáme jako lineární obal těch sloupců, které odpovídají bázickým proměnným (při Gaussovce nám tam vyjdou pivoty)
\end{itemize}
\end{itemize}
\item Isomorfismus prostorů
\begin{itemize}
\item vektorové prostory jsou isomorfní, pokud mezi nimi existuje isomorfismus, tedy bijektivní (vzájemně jednoznačné) lineární zobrazení
\item pro isomorfismus $f$ platí, že existuje $f^{-1}$ a je také isomorfismem
\item isomorfní prostory mají shodné dimenze
\item věta: $f$ je isomorfismus, právě když $[f]_{X,Y}$ je regulární
\end{itemize}
\item Skalární součin, norma indukovaná skalárním součinem
\begin{itemize}
\item skalární součin pro vektorové prostory nad komplexními čísly
\begin{itemize}
\item skalární součin na vektorovém prostoru $V$ nad $\mathbb C$ je zobrazení, které přiřadí každé dvojici vektorů $u,v\in V$ skalár $\langle u|v\rangle\in\mathbb C$ tak, že jsou splněny následující axiomy:
\begin{itemize}
\item $\forall u\in V:\langle u|u\rangle\in\mathbb R^+_0$ (reálné číslo $\geq 0$)
\item $\forall u\in V:\langle u|u\rangle=0\iff u=0$
\item $\forall u,v\in V: \langle v|u\rangle=\overline{\langle u|v\rangle}$ (komplexně sdružené)
\item $\forall u,v,w \in V: \langle u+v|w\rangle=\langle u|w\rangle+\langle v|w\rangle$
\item $\forall u,v\in V,\,\forall \alpha\in\mathbb C: \langle \alpha u|v\rangle=\alpha\langle u|v\rangle$
\end{itemize}
\item formálně je každý skalární součin zobrazení $V\times V\to\mathbb C$
\item skalární součin na $V$ nad $\mathbb R$ je zobrazení $V\times V\to\mathbb R$, přičemž $\langle v|u\rangle = \langle u|v\rangle$ (ve třetím axiomu), $\alpha\in\mathbb R$ (v posledním axiomu)
\item standardní skalární součin na $\mathbb R^n$: $\langle u|v\rangle=v^Tu$
\item standardní skalární součin na $\mathbb C^n$: $\langle u|v\rangle=v^Hu$
\end{itemize}
\item norma spojená se skalárním součinem
\begin{itemize}
\item je-li $V$ prostor se skalárním součinem, pak norma odvozená ze skalárního součinu je zobrazení $V\to R^+_0$ přiřazující vektoru $u$ jeho normu $\|u\|=\sqrt{\langle u|u\rangle}$
\item geometrická interpretace v euklidovském prostoru $\mathbb R^n$
\begin{itemize}
\item $\|u\|$ \dots{} délka (velikost) $u$
\item $\|u-v\|$ \dots{} vzdálenost bodů $u,v$
\item $\langle u|v\rangle$ souvisí s \quotedblbase{}úhlem\textquotedblleft{} $\varphi$ mezi $u,v$ a délkami $u,v$
\begin{itemize}
\item $\langle u|v\rangle=\|u\|\cdot\|v\|\cos\varphi$
\item to vyplývá z kosinové věty, kde $a=\|u\|,\,b=\|v\|,\,c=\|u-v\|$
\end{itemize}
\end{itemize}
\end{itemize}
\item kolmé vektory
\begin{itemize}
\item vektory $u,v$ z prostoru se skalárním součinem jsou kolmé (ortogonální), pokud $\langle u|v\rangle=0$
\item kolmé vektory značíme $u\perp v$
\end{itemize}
\end{itemize}
\item Pythagorova věta, Cauchyho-Schwarzova nerovnost, trojúhelníková nerovnost
\begin{itemize}
\item Pythagoras: $\|a\|^2+\|b\|^2=\|c\|^2$
\begin{itemize}
\item pokud $a,b$ jsou na sebe kolmé a $c=b-a$
\item ovšem taky platí $\|a+b\|^2=\|a\|^2+\|b\|^2$, rovněž pro kolmé $a,b$
\end{itemize}
\item Cauchy-Schwarz: $|\langle u|v\rangle|\leq\sqrt{\langle u|u\rangle\langle v|v\rangle}$
\item trojúhelníková nerovnost: $\|u+v\|\leq\|u\|+\|v\|$
\end{itemize}
\item Ortonormální systémy vektorů, Fourierovy koeficienty, Gramova-Schmidtova ortogonalizace
\begin{itemize}
\item ortonormální báze
\begin{itemize}
\item bázi $Z=\lbrace v_1,\dots,v_n \rbrace$ prostoru $V$ se skalárním součinem nazveme ortonormální, pokud platí $v_i\perp v_j$ pro každé $i\neq j$ a také $\|v_i\|=1$ pro každé $v_i\in Z$
\item pozorování: matice, jejichž sloupce tvoří vektory ortonormální báze $\mathbb C^n$ vzhledem ke std. skal. součinu, splňují $A^HA=I_n\implies$ jsou unitární
\end{itemize}
\item Fourierovy koeficienty
\begin{itemize}
\item nechť $Z=\lbrace v_1,\dots,v_n\rbrace$ je ortonormální báze prostoru $V$
\item tvrzení: pro každé $u\in V$ platí $u=\langle u|v_1\rangle v_1+\dots+\langle u|v_n\rangle v_n$
\item $\langle u|v_i\rangle$ \dots{} Fourierovy koeficienty
\end{itemize}
\item algoritmus GS ortogonalizace
\begin{itemize}
\item převede libovolnou bázi $(u_1,\dots,u_n)$ prostoru $V$ se skalárním součinem na ortonormální bázi $(v_1,\dots,v_n)$
\item for $i=1,\dots,n$ do
\begin{itemize}
\item $w_i\leftarrow u_i-\sum^{i-1}_{j=1}\langle u_i|v_j\rangle v_j$
\item $v_i\leftarrow \frac1{\|w_i\|}w_i$
\end{itemize}
\item end
\end{itemize}
\end{itemize}
\item Ortogonální doplněk, ortogonální projekce, projekce jako lineární zobrazení
\begin{itemize}
\item ortogonální doplněk
\begin{itemize}
\item ortogonální doplněk podmnožiny $V$ prostoru se skalárním součinem $W$ je $V^\perp=\lbrace u\in W:(\forall v\in V)(u\perp v)\rbrace$
\end{itemize}
\item ortogonální projekce
\begin{itemize}
\item nechť $W$ je prostor se skalárním součinem a $V$ je jeho podprostor s ortonormální bází $Z=(v_1,\dots,v_n)$
\item zobrazení $p_Z:W\to V$ definované $p_Z(u)=\sum^n_{i=1}\langle u|v_i \rangle v_i$ je ortogonální (kolmá) projekce $W$ na $V$
\end{itemize}
\item pozorování: ortogonální projekce je lineární zobrazení
\end{itemize}
\item Ortogonální matice a jejich vlastnosti
\begin{itemize}
\item matice $Q\in\mathbb R^{n\times n}$ je ortogonální, pokud $Q^TQ=I_n$
\begin{itemize}
\item pro ortogonální matici $Q$ platí $Q^T=Q^{-1}$
\end{itemize}
\item věta: $Q$ je ortogonální, právě když sloupce tvoří ortonormální bázi $\mathbb R^n$
\item zjevně: je-li $Q$ ortogonální, pak\dots{}
\begin{itemize}
\item $Q^T$ je rovněž ortogonální
\item $Q^{-1}$ existuje a je ortogonální
\end{itemize}
\item součin ortogonálních matic je ortogonální matice
\item další vlastnosti
\begin{itemize}
\item $\braket{Qx|Qy}=\braket{x|y}$
\item $\|Qx\|=\|x\|$
\item zobrazení $x\mapsto Qx$ zachovává úhly a délky
\item platí to i naopak - matice zobrazení zachovávajícího skalární součin je ortogonální
\item $\forall i,j:|Q_{ij}|\leq 1$
\item $\begin{pmatrix}1 & o^T \\ o & Q\end{pmatrix}$ je ortogonální matice
\end{itemize}
\end{itemize}
\item Definice a základní vlastnosti determinantu (multiplikativnost, determinant transponované matice, vztah s regularitou a vlastními čísly)
\begin{itemize}
\item determinant matice $A\in\mathbb K^{n\times n}$ je dán výrazem \[\text{det }A=\sum_{p\in S_n}\text{sgn}(p)\prod_{i=1}^na_{i,p(i)}\]
\item věta: $\forall A,B\in\mathbb K^{n\times n}:\text{det}(AB)=\text{det }A\cdot\text{det }B$
\item transpozice nemění determinant, $\mathrm{det}\,A=\mathrm{det}\,A^T$
\item matice je regulární, právě když má nenulový determinant
\item pro regulární matici $A$ platí $\mathrm{det}(A^{-1})=\mathrm{det}(A)^{-1}$
\item výpočet determinantu je nutné nastudovat
\begin{practicebox}
Vlastní výpočet determinantu (rozvoj, úprava na trojúhelníkový tvar) si natrénuj na příkladech.
\end{practicebox}
\item $\mathrm{det}(A)=\lambda_1\cdot\lambda_2\cdot\ldots\cdot\lambda_n$
\item $\lambda$ je vlastním číslem $A$, právě když $\mathrm{det}(A-\lambda I_n)=0$
\begin{itemize}
\item tedy právě když je kořenem charakteristického polynomu $p_A(t)=\mathrm{det}(A-tI_n)$
\end{itemize}
\end{itemize}
\item Laplaceův rozvoj determinantu
\begin{itemize}
\item notace: $A^{ij}$ je podmatice získaná z $A$ odstraněním $i$-tého řádku a $j$-tého sloupce
\item věta: pro libovolné $A\in\mathbb K^{n\times n}$ a jakékoli $i\in \lbrace 1,\dots,n\rbrace$ platí \[\text{det }A=\sum^n_{j=1}a_{ij}(-1)^{i+j}\text{ det }A^{ij}\]
\end{itemize}
\item Geometrická interpretace determinantu
\begin{itemize}
\item objem rovnoběžnostěnu určeného $k$ vektory je roven absolutní hodnotě determinantu matice, jejíž sloupce tyto vektory tvoří
\item po provedení lineárního zobrazení se objemy těles změní úměrně absolutní hodnotě determinantu matice zobrazení
\end{itemize}
\item Definice, geometrický význam a základní vlastnosti vlastních čísel, charakteristický polynom, násobnost vlastních čísel
\begin{itemize}
\item definice
\begin{itemize}
\item mějme matici $A\in\mathbb C^{n\times n}$
\item vlastní číslo matice $A$ je jakékoliv $\lambda\in\mathbb C$, pro které existuje vektor $x\in\mathbb C^n\setminus \lbrace o\rbrace$ takový, že $Ax=\lambda x$
\item $x$ je pak vlastní vektor příslušný vlastnímu číslu $\lambda$
\end{itemize}
\item geometrie
\begin{itemize}
\item vlastní vektor reprezentuje invariantní směr při zobrazení $x\mapsto Ax$, tedy směr, který se zobrazí opět na ten samý směr
\item vlastní číslo odpovídá škálování v tomto invariantním směru
\end{itemize}
\item základní vlastnosti
\begin{itemize}
\item determinant matice je roven součinu vlastních čísel
\item stopa matice (součet diagonály) je rovna součtu vlastních čísel
\item matice je regulární, právě když nula není její vlastní číslo
\item $A^T$ má stejná vlastní čísla jako $A$, ale vlastní vektory obecně jiné
\item uvažujeme matici $A\in\mathbb C^{n\times n}$ s vlastními čísly $\lambda_1,\dots,\lambda_n$ a jim odpovídajícími vlastními vektory $x_1,\dots,x_n$
\begin{itemize}
\item je-li $A$ regulární, pak $A^{-1}$ má vlastní čísla $\lambda_1^{-1},\dots,\lambda_n^{-1}$ a vlastní vektory $x_1,\dots,x_n$
\item $A^2$ má vlastní čísla $\lambda_1^{2},\dots,\lambda_n^{2}$ a vlastní vektory $x_1,\dots,x_n$
\item podobně pro $\alpha A$ a taky pro $A+\alpha I_n$
\end{itemize}
\end{itemize}
\item charakteristický polynom \dots{} $p_A(t)=\mathrm{det}(A-tI_n)$
\begin{itemize}
\item vlastní čísla matice $A$ jsou právě kořeny jejího charakteristického polynomu a je jich $n$ včetně násobností
\end{itemize}
\item algebraická násobnost $\lambda$ je rovna násobnosti $\lambda$ jakožto kořene $p_A$
\item geometrická násobnost $\lambda$ je rovna $n-\mathrm{rank}(A-\lambda I_n)$, tj. počtu lineárně nezávislých vlastních vektorů, které odpovídají $\lambda$
\item algebraická násobnost $\geq$ geometrická násobnost
\begin{itemize}
\item obvykle se násobností myslí ta algebraická
\end{itemize}
\end{itemize}
\item Podobnost a diagonalizovatelnost matic, spektrální rozklad
\begin{itemize}
\item matice $A,B\in\mathbb C^{n\times n}$ jsou podobné, pokud existuje regulární $S\in\mathbb C^{n\times n}$ tak, že $A=SBS^{-1}$
\item věta: podobné matice mají stejná vlastní čísla
\begin{itemize}
\item počet vlastních vektorů pro konkrétní vlastní číslo je stejný u obou matic
\end{itemize}
\item matice $A\in\mathbb C^{n\times n}$ je diagonalizovatelná, je-li podobná nějaké diagonální matici
\item diagonalizovatelná matice $A$ jde tedy vyjádřit ve tvaru $A=S\Lambda S^{-1}$, kde $S$ je regulární a $\Lambda$ diagonální
\begin{itemize}
\item tomuto se říká spektrální rozklad
\item sloupce matice $S$ odpovídají vlastním vektorům (v~pořadí, v němž jsou vlastní čísla na diagonále $\Lambda$)
\end{itemize}
\item věta: matice $A\in\mathbb C^{n\times n}$ je diagonalizovatelná, právě když má $n$ lineárně nezávislých vlastních vektorů
\item věta: pokud má matice $n$ navzájem různých vlastních čísel, pak je diagonalizovatelná
\item $A^2=S\Lambda^2 S^{-1}$
\end{itemize}
\item Symetrické matice, jejich vlastní čísla a spektrální rozklad
\begin{itemize}
\item matice je symetrická, pokud $A=A^T$
\item vlastní čísla reálných symetrických matic jsou reálná
\item věta: pro každou symetrickou matici $A\in\mathbb R^{n\times n}$ existuje ortogonální $Q\in\mathbb R^{n\times n}$ a diagonální $\Lambda\in\mathbb R^{n\times n}$ takové, že $A=Q\Lambda Q^T$
\end{itemize}
\item Positivně semidefinitní a positivně definitní matice - charakterizace a vlastnosti, vztah se skalárním součinem, vlastními čísly
\begin{itemize}
\item definice
\begin{itemize}
\item buď $A\in\mathbb R^{n\times n}$ symetrická
\item pak $A$ je positivně semidefinitní, pokud $x^TAx\geq 0$ pro všechna $x\in\mathbb R^n$
\item $A$ je positivně definitní, pokud $x^TAx\gt 0$ pro všechna $x\neq o$
\end{itemize}
\item poznámka: nesymetrické matice můžeme symetrizovat úpravou $\frac12(A+A^T)$
\item následující tří vlastnosti jsou ekvivalentní
\begin{itemize}
\item $A$ je positivně semidefinitní (nebo definitní)
\item vlastní čísla $A$ jsou nezáporná
\begin{itemize}
\item pro definitní musí být kladná
\end{itemize}
\item existuje matice $U\in\mathbb R^{m\times n}$ taková, že $A=U^TU$
\begin{itemize}
\item pro definitní musí mít matice $U$ hodnost $n$
\end{itemize}
\end{itemize}
\item vlastnosti
\begin{itemize}
\item jsou-li $A,B$ positivně definitní, pak i $A+B$ je p. d.
\item je-li $A$ positivně definitní a $\alpha\gt 0$, pak i $\alpha A$ je p. d.
\item je-li $A$ positivně definitní, pak je regulární a $A^{-1}$ je p. d.
\end{itemize}
\item operace $\braket{x|y}$ je skalárním součinem v $\mathbb R^n$, právě když má tvar $\braket{x|y}=x^TAy$ pro nějakou positivně definitní matici $A\in\mathbb R^{n\times n}$
\item testování p. d. pomocí Gaussovy eliminace
\begin{itemize}
\item používáme pouze úpravu přičtení násobku řádku s pivotem k jinému řádku pod ním
\item pokud nám vyjde odstupňovaný tvar s kladnou diagonálou, byla původní matice positivně definitní
\end{itemize}
\end{itemize}
\item Choleského rozklad (znění věty a praktické použití)
\begin{itemize}
\item pro každou positivně definitní matici $A\in\mathbb R^{n\times n}$ existuje jediná dolní trojúhelníková matice $L\in\mathbb R^{n\times n}$ s kladnou diagonálou taková, že $A=LL^T$
\item algoritmus vypadá složitě, ale dá se to vymyslet intuitivně (je vhodné začít tím, že první prvek matice $A$ je druhou mocninou čísla, které je v obou trojúhelníkových maticích vlevo nahoře)
\item praktické použití
\begin{itemize}
\item řešíme soustavu $Ax=b$
\item máme Choleského rozklad $A=LL^T$
\item snadno najdeme řešení soustavy $Ly=b$ (díky nulám v trojúhelníkové matici stačí substituovat)
\item pak najdeme řešení soustavy $L^Tx=y$
\item proč to funguje?
\begin{itemize}
\item $L(\underbrace{L^Tx}_{=y})=b$
\end{itemize}
\end{itemize}
\end{itemize}
\end{itemize}
\subsection{3. Diskrétní matematika}
\begin{itemize}
\item Relace, vlastnosti binárních relací (reflexivita, symetrie, antisymetrie, tranzitivita)
\begin{itemize}
\item (binární) relace mezi množinami $X,Y$ je podmnožina $X\times Y$
\item relace na množině $X$ je podmnožina $X^2$
\item značení - pro relaci $R$ mezi $X,Y:xRy \equiv (x,y)\in R$
\item příklady relací
\begin{itemize}
\item prázdná $\emptyset$
\item univerzální $X\times Y$
\item diagonální $\Delta_X := \lbrace (x,x)\mid x\in X\rbrace$, např. rovnost $x=y$
\end{itemize}
\item operace s relacemi
\begin{itemize}
\item inverze
\begin{itemize}
\item k relaci $R$ mezi $X,Y$ lze definovat inverzní relaci $R^{-1}$ mezi $Y,X$, přičemž $R^{-1} := \lbrace(y,x)\mid (x,y)\in R\rbrace$
\end{itemize}
\item skládání
\begin{itemize}
\item pro relaci $R$ mezi $X,Y$ a relaci $S$ mezi $Y,Z$ lze definovat složenou relaci $T=R\circ S$ mezi $X,Z$
\item $xTz \equiv \exists y \in Y: xRy \land ySz$
\item $R\circ \Delta_Y =R,\quad \Delta_X\circ R = R$
\item značení skládání funkcí: $(f\circ g)(x)=g(f(x))$
\end{itemize}
\end{itemize}
\item relace $R$ na $X$ je\dots{}
\begin{itemize}
\item reflexivní $\equiv \forall x \in X: xRx$
\begin{itemize}
\item $\Delta_X \subseteq R$
\end{itemize}
\item symetrická $\equiv \forall x,y \in X: xRy \implies yRx$
\begin{itemize}
\item $R=R^{-1}$
\end{itemize}
\item antisymetrická $\equiv \forall x,y \in X: xRy \land yRx \implies x=y$
\begin{itemize}
\item $R\cap R^{-1}\subseteq \Delta_X$
\end{itemize}
\item tranzitivní $\equiv \forall x,y,z \in X: xRy \land yRz \implies xRz$
\begin{itemize}
\item $R\circ R \subseteq R$
\end{itemize}
\end{itemize}
\end{itemize}
\item Ekvivalence a rozkladové třídy
\begin{itemize}
\item relace $R$ na $X$ je ekvivalence $\equiv R$ je reflexivní \& symetrická \& tranzitivní
\begin{itemize}
\item např. rovnost čísel, rovnost mod $K$, geometrická podobnost
\end{itemize}
\item ekvivalenční třída prvku $x \in X:R[x]=\lbrace y \in X : xRy\rbrace$
\item množinový systém $\mathcal S\subseteq 2^X$ je rozklad množiny $X \equiv$
\begin{itemize}
\item $\forall A \in \mathcal S: A \neq \emptyset$
\item $\forall A,B\in \mathcal S: A\neq B \implies A \cap B = \emptyset$
\item $\bigcup_{A\in \mathcal S}A=X$
\end{itemize}
\item věta (vztah mezi ekvivalencemi a rozklady)
\begin{enumerate}
\item $\forall x \in X: R[x]\neq \emptyset$
\item $\forall x,y \in X:$ buď $R[x]=R[y]$, nebo $R[x] \cap R[y]=\emptyset$
\item $\lbrace R[x] \mid x \in X\rbrace$ (množina všech ekvivalenčních tříd) určuje ekvivalenci $R$ jednoznačně
\end{enumerate}
\end{itemize}
\item Částečná uspořádání - základní pojmy (minimální a maximální prvky, nejmenší a největší prvky, řetězec, antiřetězec)
\begin{itemize}
\item relace $R$ na množině $X$ je (částečné) uspořádání $\equiv R$ je reflexivní \& antisymetrická \& tranzitivní
\item (částečně) uspořádaná množina $(X,R)$
\begin{itemize}
\item zkráceně ČUM
\item $R$ je (částečné) uspořádání na $X$
\end{itemize}
\item prvky $x,y \in X$ jsou porovnatelné $\equiv xRy \lor yRx$
\item uspořádání je lineární $\equiv \forall x,y\in X$ porovnatelné
\begin{itemize}
\item (všechny prvky množiny jsou navzájem porovnatelné)
\end{itemize}
\item částečné uspořádání (nebo pouze uspořádání) je obecný pojem, některá taková uspořádání jsou navíc lineární
\item ostré uspořádání - každému uspořádání $\leq$ na X přiřadíme relaci $<$ na X: $a < b \equiv a \leq b \land a \neq b$ 
\begin{itemize}
\item pozor - ostré uspořádání není speciálním případem uspořádání (protože není reflexivní)
\item vlastnosti ostrého uspořádání - ireflexivní, antisymetrické, tranzitivní
\end{itemize}
\item Hasseův diagram graficky zachycuje vztahy mezi prvky ČUM (porovnatelné prvky jsou spojeny, větší prvky jsou výše)
\item prvek $x\in X$ je nejmenší $\equiv \forall y \in X: x\leq y$
\item prvek $x\in X$ je minimální $\equiv \nexists y \in X: y\lt x$
\item $x$ je nejmenší $\implies$ $x$ je minimální
\item v Hassově diagramu
\begin{itemize}
\item z minimálního prvku dolů nevede žádná spojnice
\item nejmenší prvek je nejníž v diagramu, existuje do něj cesta z libovolného jiného prvku
\end{itemize}
\item věta: konečná neprázdná uspořádaná množina má minimální a maximální prvek
\item pro $(X,\leq)$ ČUM:
\begin{itemize}
\item $A\subseteq X$ je řetězec $\equiv \forall a,b \in A: a,b$ jsou porovnatelné
\item $A \subseteq X$ je antiřetězec (nezávislá množina) $\equiv \nexists a,b$ různé \& porovnatelné
\end{itemize}
\end{itemize}
\item Výška a šířka částečně uspořádané množiny a věta o jejich vztahu (o dlouhém a širokém)
\begin{itemize}
\item $\omega(X,\leq)$ je délka nejdelšího řetězce = maximum z délek řetězců (výška uspořádání)
\item $\alpha(X,\leq)$ je délka nejdelšího antiřetězce (šířka uspořádání)
\item věta: pro každou konečnou ČUM $(X,\leq)$ platí $\alpha(X,\leq)\cdot \omega(X,\leq)\geq |X|$
\item důsledek: $\text{max}(\alpha,\omega)\geq\sqrt{|X|}$
\end{itemize}
\item Funkce, typy funkcí (prostá, na, bijekce)
\begin{itemize}
\item funkce z množiny $X$ do množiny $Y$ je relace $A$ mezi $X$ a $Y$ t. ž. $(\forall x\in X)(\exists! y\in Y):xAy$
\item funkce $f:X\to Y$ je\dots{}
\begin{itemize}
\item prostá (injektivní) $\equiv \nexists x,x' \in X: x \neq x' \land f(x)=f(x')$
\item na $Y$ (surjektivní) $\equiv (\forall y \in Y)(\exists x \in X): f(x)=y$
\item vzájemně jednoznačná (bijektivní) $\equiv (\forall y \in Y)(\exists! x \in X):f(x)=y$
\begin{itemize}
\item taková funkce je tedy prostá i \quotedblbase{}na\textquotedblleft{}
\item k takové funkci existuje inverzní funkce $f^{-1}$ z $Y$ do $X$
\end{itemize}
\end{itemize}
\end{itemize}
\item Počty různých typů funkcí mezi dvěma konečnými množinami
\begin{itemize}
\item věta: počet $f:N\to M=m^n$
\begin{itemize}
\item pro $|N|=n, |M|=m;\quad m,n\gt 0$
\end{itemize}
\item definice: klesající mocnina
\begin{itemize}
\item $m^{\underline n}=\underbrace{m\cdot (m-1)\cdot (m-2)\cdot\ldots\cdot(m-n+1)}_n$
\item tedy $m^{\underline n}=\frac{m!}{(m-n)!}$
\end{itemize}
\item věta: počet prostých $f:N\to M=m^{\underline{n}}$
\item počet surjektivních zobrazení lze určit pomocí principu inkluze a exkluze
\item bijekcí je $n!$
\item počet uspořádaných $k$-tic $|X^k|=|X|^k$, lze jej totiž vyjádřit jako počet funkcí $f:[k]\to X$
\end{itemize}
\item Permutace a jejich základní vlastnosti (počet a pevný bod)
\begin{itemize}
\item definice: $[n]=\lbrace1,2,\dots,n\rbrace$
\item věta: na množině $[n]$ existuje $n!$ permutací (podobně na každé n-prvkové množině)
\item pevný bod permutace je prvek, který se zobrazí sám na sebe
\end{itemize}
\item Kombinační čísla a vztahy mezi nimi, binomická věta a její aplikace
\begin{itemize}
\item kombinační číslo (binomický koeficient) ${n\choose k}:={n^{\underline k} \over k!}={n!\over k!(n-k)!}$
\item Pascalův trojúhelník - n roste shora dolů, k zleva doprava
\item $X\choose k$ \dots{} množina všech k-prvkových podmnožin množiny $X$
\item ${n\choose k} = {n\choose n-k}$ \dots{} každé k-prvkové podmnožině přiřadíme její doplněk
\begin{itemize}
\item Pascalův trojúhelník je symetrický
\end{itemize}
\item ${n-1\choose k-1}+{n-1\choose k}={n\choose k}$
\begin{itemize}
\item zvolíme jeden prvek $a$ a rozdělíme všechny k-prvkové podmnožiny podle toho, zda obsahují $a$, nebo ne
\item buňka v Pascalově trojúhelníku je součtem dvou buněk nad ní (pokud není na kraji)
\end{itemize}
\item Binomická věta: $(x+y)^n=\sum_{i=0}^n{n\choose i}x^{n-i}y^i$
\item aplikace
\begin{itemize}
\item řádek Pascalova trojúhelníku se sečte na $2^n=(1+1)^n=\sum_{k=0}^n\binom{n}{k}$
\item důkaz $(x^n)'=nx^{n-1}$
\begin{itemize}
\item použijeme $f(x+h)=(x+h)^n=x^n+nx^{n-1}h+\binom n2 x^{n-2}h^2+\dots+h^n$
\item dosadíme do $\lim_{h\to 0}\frac{f(x+h)-f(x)}{h}$
\end{itemize}
\item $e=\lim\,(1+\frac1n)^n$
\end{itemize}
\end{itemize}
\item Princip inkluze a exkluze: obecná formulace (a důkaz)
\begin{itemize}
\item konkrétní ukázky principu inkluze a exkluze (PIE)
\begin{itemize}
\item $|A\cup B|=|A|+|B|-|A\cap B|$
\item $|A\cup B\cup C|=|A|+|B|+|C|-|A\cap B|-|A\cap C|-|B\cap C|+|A\cap B\cap C|$
\end{itemize}
\item věta: pro konečné množiny $A_1,\dots,A_n$ platí \[|\bigcup_{i=1}^nA_i|=\sum_{k=1}^n(-1)^{k+1}\sum_{I\in{[n]\choose k}}|\bigcap_{i\in I}A_i|\]
\item důkaz
\begin{itemize}
\item pro prvek $x$ ve sjednocení spočítáme příspěvky k levé (vždy 1) a pravé straně
\item nechť $x$ patří do právě $t$ množin
\item průniky $k$-tic
\begin{itemize}
\item $k \gt t$ \dots{} přispěje $0$
\item $k\leq t$ \dots{} přispěje $(-1)^{k+1}{t\choose k}$
\begin{itemize}
\item vybíráme $k$-tice množin z $t$-množin, do kterých prvek patří
\item minus jednička vychází ze vzorce
\end{itemize}
\item chceme $\sum_{k=1}^t(-1)^{k+1}{t\choose k}=1$
\item lze upravit na $\sum_{k=1}^t(-1)^{k}{t\choose k}=-1$
\item z binomické věty $0=(1-1)^t=\sum_{k=0}^t{t\choose k}(-1)^{k}$
\item tedy bez prvního členu se součet rovná $-1\quad \square$
\end{itemize}
\end{itemize}
\end{itemize}
\item Princip inkluze a exkluze: použití (problém šatnářky, Eulerova funkce pro počet dělitelů, počet surjekcí)
\begin{itemize}
\item problém šatnářky
\begin{itemize}
\item Šatnářka $n$ pánům vydá náhodně $n$ klobouků (které si předtím odložili v~šatně). Jaká je pravděpodobnost, že žádný pán nedostane od šatnářky zpět svůj klobouk?
\item jaká je pravděpodobnost, že náhodně zvolená permutace nebude mít žádný pevný bod
\begin{itemize}
\item každá z $n!$ permutací je stejně pravděpodobná
\item $\text{š}(n)$ \dots{} počet permutací bez pevného bodu
\item pravděpodobnost je rovna $\text{š}(n)/n!$
\end{itemize}
\item $S_n$ \dots{} množina všech permutací
\item $A_i=\lbrace \pi \in S_n : \pi(i)=i\rbrace$
\begin{itemize}
\item množina permutací s pevným bodem v $i$-tém prvku (jedna permutace může patřit do více takových množin)
\end{itemize}
\item $A=\bigcup_{i=1}^n A_i$ \dots{} (množina všech \quotedblbase{}špatných\textquotedblleft{} permutací)
\item musíme vyjádřit velikosti průniků
\begin{itemize}
\item permutací s (konkrétními) $k$ pevnými body je $(n-k)!$, protože permutuji všechny prvky kromě těch pevných
\end{itemize}
\item dosazením do principu inkluze a exkluze vyjde $|A|=\sum_{k=1}^n(-1)^{k+1}{n\choose k}(n-k)!$
\begin{itemize}
\item $n\choose k$ vyplývá z počtu prvků druhé sumy
\item ${n\choose k}(n-k)!={n!\over k!}$
\end{itemize}
\item $|A|=\sum_{k=1}^n(-1)^{k+1}\frac{n!}{k!}=n!\cdot\sum_{k=1}^n{(-1)^{k+1}\over k!}$
\item $|A|=n!({1\over 1!}-{1\over 2!}+{1\over 3!}-\dots+{(-1)^{n+1}\over n!})$
\item $\text{š}(n)=n!-|A|=n!\cdot(1-{1\over 1!}+{1\over 2!}-{1\over 3!}+\dots+{(-1)^n\over n!})$
\begin{itemize}
\item závorka konverguje k $e^{-1}$
\item závorka odpovídá pravděpodobnosti v problému šatnářky
\end{itemize}
\item $\text{š}(n)=n!\cdot\sum_{k=0}^n\frac{(-1)^k}{k!}$
\item pravděpodobnost \dots{} $\sum_{k=0}^n\frac{(-1)^k}{k!}\approx e^{-1}$
\end{itemize}
\item Eulerova funkce $\varphi(n)$
\begin{itemize}
\item je to počet čísel z $[n]$ nesoudělných s $n$
\begin{itemize}
\item pro prvočíslo $p$ platí $\varphi(n)=p-1$, protože je soudělné samo se sebou a protože podle definice nesoudělnosti (NSD = 1) není soudělné s~jedničkou
\end{itemize}
\item tvrzení: nechť $n=p_1^{e_1}\cdot p_2^{e_2}\cdot\ldots\cdot p_k^{e_k}$ je prvočíselný rozklad, pak $\varphi(n)=n(1-\frac{1}{p_1})(1-\frac1{p_2})\dots(1-\frac1{p_k})$
\item důkaz
\begin{itemize}
\item jako $A_i$ označím množinu čísel z $[n]$ soudělných $i$-tým prvočíslem $p_i$
\item zjevně $\varphi(n)=n-|\bigcup_{i=1}^k A_i|$
\item $|A_i|=\frac n{p_i}$
\item pro různé $i,j\in[k]$ platí $|A_i\cap A_j|=\frac n{p_ip_j}$
\item to se dá zobecnit na průnik více množin
\item dle PIE $|\bigcup_{i=1}^kA_i|=\sum_{I\in 2^{[n]}\setminus\set\emptyset}(-1)^{|I|+1}\frac{n}{\prod_{i\in I} p_i}$
\item $=n(\frac1{p_1}+\dots+\frac1{p_k}-\frac1{p_1p_2}-\dots+\frac1{p_1p_2p_3}+\dots+(-1)^{k+1}\frac1{p_1\dots p_k})$
\item proto $\varphi(n)=n(1-\frac1{p_1}-\dots-\frac1{p_k}+\frac{1}{p_1p_2}+\dots+(-1)^k\frac1{p_1\dots p_k})$
\item tedy $\varphi(n)=n(1-\frac{1}{p_1})(1-\frac1{p_2})\dots(1-\frac1{p_k})\quad\square$
\begin{itemize}
\item z každé závorky vždy vybíráme levý nebo pravý člen
\end{itemize}
\end{itemize}
\end{itemize}
\item počet surjekcí
\begin{itemize}
\item věta
\begin{itemize}
\item nechť $X,Y$ jsou konečné množiny, $|X|=n$, $|Y|=m$, $n\geq m\gt 0$
\item potom počet zobrazení $f$ množiny $X$ na množinu $Y$ je $\sum_{k=0}^m(-1)^k\binom mk(m-k)^n$
\end{itemize}
\item důkaz
\begin{itemize}
\item počet všech funkcí z $X$ do $Y$ je $m^n$
\item definujeme $F_i$ jako množinu zobrazení, že $i$-tý prvek z množiny $Y$ není pokrytý (neexistuje $x\in X$, aby $f(x)=y_i$, kde $y_i$ je $i$-tý prvek z $Y$)
\item takže počet surjekcí bude $m^n-|\bigcup_{i=1}^m F_i|$
\item použijeme PIE: $|\bigcup_{i=1}^mF_i|=\sum_{k=1}^m(-1)^{k+1}\sum_{I\in{[m]\choose k}}|\bigcap_{i\in I}F_i|$
\item $=\sum_{k=1}^m(-1)^{k+1}\binom mk(m-k)^n$
\begin{itemize}
\item $(m-k)^n$ dostaneme tak, že uvažujeme funkce z $n$-prvkové množiny do množiny s $m-k$ prvky (na vybraných $k$ bodů naše funkce nesmí směřovat)
\end{itemize}
\item počet surjekcí $m^n-|\bigcup_{i=1}^m F_i|=m^n-\sum_{k=1}^m(-1)^{k+1}\binom mk(m-k)^n$ pak ještě upravíme do finálního tvaru $\quad\square$
\end{itemize}
\end{itemize}
\end{itemize}
\item Hallova věta o systému různých reprezentantů a její vztah k párování v bipartitním grafu, princip důkazu a algoritmické aspekty (polynomiální algoritmus pro nalezení SRR)
\begin{itemize}
\item párování je taková množina hran, kde žádné dvě hrany nemají společný vrchol
\begin{itemize}
\item maximální párování je takové párování, které má nejvíce hran
\item perfektní párování je takové párování, které pokrývá všechny vrcholy grafu
\end{itemize}
\item systém různých reprezentantů v hypergrafu
\begin{itemize}
\item systém různých reprezentantů (SRR) v hypergrafu $H=(V,E)$ je funkce $r:E\to V$ taková, že
\begin{itemize}
\item $\forall e\in E: r(e)\in e$
\item $\forall e,f\in E:e\neq f\implies r(e)\neq r(f)$ \dots{} tj. $r$ je prostá
\end{itemize}
\item $r(e)$ \dots{} \quotedblbase{}reprezentant hyperhrany $e$\textquotedblleft{}
\item analogie s předsedy spolků
\end{itemize}
\item pozorování: v bipartitním grafu s partitami $A,B$ má každé párování velikost nejvýše $|A|$ (i nejvýše $|B|$)
\item definice: pro graf $G=(V,E)$ a množinu $X\subseteq V$ označím $N(X):=\set{y\in V\setminus X: (\exists x\in X)\set{x,y}\in E}$
\item Hallova věta, bipartitní grafová verze
\begin{itemize}
\item nechť $G$ je bipartitní graf s partitami $A,B$, potom $G$ má párování velikosti $|A|$ $\iff\underbrace{\forall X\subseteq A:|N(X)|\geq |X|}_{\text{\quotedblbase{}Hallova podm\text{í}nka\textquotedblleft{}}}$
\end{itemize}
\item důkaz
\begin{itemize}
\item $\implies$
\begin{itemize}
\item pokud existuje párování velikosti $|A|$, tak pro každou $X\subseteq A$ existuje $|X|$ vrcholů spárovaných s $X$, ty patří do $N(X)$, tedy $|N(X)|\geq |X|$
\end{itemize}
\item $\impliedby$
\begin{itemize}
\item nechť $M$ je největší párování v $G$
\item pro spor nechť $|M|\lt|A|$
\item dle K\H{o}nigovy-Egerváryho věty existuje pokrytí $C$, kde $|C|=|M|\lt |A|$
\item $C_A:= C\cap A$
\item $C_B:=C\cap B$
\item $X:=A\setminus C_A$
\item zjistím, že $N(X)\subseteq C_B$
\begin{itemize}
\item protože nepokryté hrany musí pokrývat vrcholy v $N(X)$
\end{itemize}
\item navíc $|X|=|A|-|C_A|\gt |C_B|\geq |N(X)|$
\begin{itemize}
\item jelikož $|C|\lt|A|\implies |A|\gt|C_A|+|C_B|$
\end{itemize}
\item to je spor s Hallovou podmínkou
\end{itemize}
\item idea $\impliedby$
\begin{itemize}
\item uvažujeme ostře menší párování, tedy podle K\H{o}nigovy-Egerváryho věty musí existovat i menší pokrytí
\begin{itemize}
\item věta říká, že velikost maximální párování se rovná velikosti minimálního pokrytí (jako u toku a řezu)
\end{itemize}
\item prohlásíme, že vrcholy v $X$ nejsou v pokrytí, tak ty hrany mezi nimi a $N(X)$ musí pokrývat vrcholy v $N(X)$, ale těch by pak bylo málo (méně než požaduje Hallova podmínka)
\item takže jsme dokázali obměnu implikace
\end{itemize}
\end{itemize}
\item pozorování: jakmile uvnitř hypergrafu je množina čtyř hyperhran, které ve svém sjednocení mají tři vrcholy, tak nemůžu najít systém různých reprezentantů
\begin{itemize}
\item platí to obecně pro množinu $k$ hyperhran - když budou mít ve sjednocení méně než $k$ vrcholů, tak nenajdu SRR
\item implikace platí i obráceně, lze vyslovit jako větu
\end{itemize}
\item Hallova věta, hypergrafová verze
\begin{itemize}
\item hypergraf $H=(V,E)$ má SRR, právě když $\forall F\subseteq E:\left|\bigcup_{e\in F} e\right|\geq |F|$
\end{itemize}
\item důkaz
\begin{itemize}
\item nechť $H=(V,E)$ je hypergraf, nechť $I_H$ je jeho graf incidence
\item $H$ má SRR $\iff$ $I_H$ má párování velikosti $|E|$
\item Hallova podmínka pro $H$: $\forall F\subseteq E: \left|\bigcup_{e\in F}e\right|\geq |F|\iff$ bipartitní Hallova podmínka pro $I_H$ a partitu $E$
\end{itemize}
\item algoritmické aspekty
\begin{itemize}
\item maximální párování lze najít v polynomiálním čase
\item proto i SRR lze najít v polynomiálním čase
\end{itemize}
\end{itemize}
\end{itemize}
\subsection{4. Teorie grafů}
\begin{itemize}
\item Základní pojmy teorie grafů - graf, vrcholy a hrany, izomorfismus grafů, podgraf, okolí vrcholu a stupeň vrcholu, doplněk grafu, bipartitní graf
\begin{itemize}
\item graf je $(V, E)$, kde $V$ je konečná neprázdná množina vrcholů a $E \subseteq {V \choose 2}$ je množina hran
\begin{itemize}
\item lze značit jako $G=(V,E)$, potom $V(G)$ je množina vrcholů a $E(G)$ je množina hran
\end{itemize}
\item izomorfismus
\begin{itemize}
\item grafy jsou izomorfní $\equiv$ existuje bijekce, která zachovává vlastnost být spojen hranou
\item v podstatě stačí přejmenovat vrcholy a dostaneme dva stejné grafy
\item značení $\cong$
\item $\cong$ je ekvivalence na libovolné množině grafů
\begin{itemize}
\item neexistuje množina všech grafů (protože neexistuje množina všech množin)
\end{itemize}
\end{itemize}
\item podgraf
\begin{itemize}
\item graf $G'=(V',E')$ je podgrafem grafu $G=(V,E)$ (značíme $G'\subseteq G$), když $V'\subseteq V \land E'\subseteq E$
\item graf $G'=(V',E')$ je indukovaným podgrafem grafu $G=(V,E)$, když $V'\subseteq V \land E'= E\cap {V'\choose 2}$
\begin{itemize}
\item \quotedblbase{}podgraf indukovaný množinou vrcholů\textquotedblleft{}
\item $G[A] := (A, E(G) \cap {A \choose 2})$, kde $A\subseteq V(G)$
\end{itemize}
\end{itemize}
\item okolí vrcholu
\begin{itemize}
\item okolí vrcholu $v$ jsou ty vrcholy, které s vrcholem $v$ sousedí (sdílejí hranu)
\item $N_G(v)=\set{u\in V(G):\exists uv\in E(G)}$
\end{itemize}
\item stupeň vrcholu
\begin{itemize}
\item stupeň vrcholu - počet hran, kterých se účastní daný vrchol
\item graf je k-regulární, pokud je stupeň všech vrcholů grafu roven k
\item skóre grafu = posloupnost stupňů vrcholů (až na pořadí) $\to$ jakmile dvěma grafům vyjde jiné skóre, nemohou být izomorfní
\end{itemize}
\item doplněk grafu
\begin{itemize}
\item doplněk/komplement grafu $G$ je graf $G_0$, který má stejné vrcholy jako $G$, ale má právě ty hrany, které v grafu $G$ chybí
\end{itemize}
\item bipartitní graf
\begin{itemize}
\item graf (V, E) je bipartitní $\equiv \exists L,P\subseteq V$ t. ž.:
\begin{itemize}
\item $L \cup P = V$
\item $L \cap P = \emptyset$
\item $\forall e \in E: |e \cap L| = 1\quad (\land\ |e \cap P| = 1)$
\begin{itemize}
\item nebo $E(G)\subseteq\lbrace\lbrace x,y\rbrace\mid x\in L,y\in P\rbrace$
\end{itemize}
\end{itemize}
\item partity grafu - jednotlivé \quotedblbase{}strany\textquotedblleft{}
\end{itemize}
\end{itemize}
\item Základní příklady grafů - úplný graf a úplný bipartitní graf, cesty a kružnice
\begin{itemize}
\item úplný graf $K_n$
\begin{itemize}
\item $V(K_n):=[n]$
\item $E(K_n):={V(K_n)\choose 2}$
\end{itemize}
\item úplný bipartitní $K_{m,n}$
\begin{itemize}
\item každý prvek nalevo je spojený s každým napravo
\item prvky na jedné straně mezi sebou nejsou spojeny
\end{itemize}
\item prázdný graf $E_n$
\begin{itemize}
\item $V(E_n):=[n]$
\item $E(E_n) =\emptyset$
\end{itemize}
\item cesta $P_n$
\begin{itemize}
\item $V(P_n):=\{0,\dots, n\}$
\item $E(P_n):=\{\{i,i+1\}\mid 0\leq i \lt n\}$
\item délka cesty se měří v počtu hran
\end{itemize}
\item kružnice/cyklus $C_n$
\begin{itemize}
\item $n\geq 3$
\item $V(C_n):=\{0,\dots, n-1\}$
\item $E(C_n):=\{\{i,(i+1)\bmod n\}\mid 0\leq i \lt n\}$
\end{itemize}
\end{itemize}
\item Souvislost grafů, komponenty souvislosti, vzdálenost v grafu
\begin{itemize}
\item graf $G$ je souvislý $\equiv\forall u,v\in V(G):$ existuje cesta v $G$ s krajními vrcholy $u,v$
\item dosažitelnost v $G$ je relace $\sim$ na $V(G)$ t. ž. $u\sim v\equiv$ existuje cesta v $G$ s krajními vrcholy $u,v$
\begin{itemize}
\item relace $\sim$ je ekvivalence
\item tranzitivita se dokazuje pomocí dvou posloupností vrcholů $x \sim y$ a $y\sim z$ a následně zvolení nejzazšího vrcholu z posloupnosti $y\sim z$, který je obsažen v posloupnosti $x\sim y$ a v tomto vrcholu se posloupnosti slepí (přičemž části za ním v první posloupnosti a před ním v druhé posloupnosti se ustřihnou)
\end{itemize}
\item komponenty souvislosti jsou podgrafy indukované třídami ekvivalence $\sim$
\begin{itemize}
\item komponenty jsou souvislé
\item graf je souvislý $\iff$ má 1 komponentu
\end{itemize}
\item vzdálenost (grafová metrika) v souvislém grafu $G$
\begin{itemize}
\item $d_G:V^2\to\mathbb R$
\item $d_G(u,v):=$ min. z délek (počtu hran) všech cest mezi $u,v$
\end{itemize}
\item vlastnosti metriky (funkce je metrika = chová se jako vzdálenost)
\begin{itemize}
\item $d_G(u,v)\geq 0$
\item $d_G(u,v) = 0 \iff u=v$
\item platí trojúhelníková nerovnost $d_G(u,w) \leq d_G(u,w) + d_G(w,v)$
\item $d_G(v,u)=d_G(u,v)$
\end{itemize}
\item grafové operace: přidání/odebrání vrcholu/hrany, dělení hrany, kontrakce hrany
\item sled v grafu (walk) - můžou se opakovat vrcholy i hrany
\begin{itemize}
\item $(v_0,e_1,v_1,e_2,\dots,e_n,v_n)$
\item $\forall i: e_i=\lbrace v_{i-1},v_i\rbrace$
\end{itemize}
\item tah v grafu - můžou se opakovat vrcholy, hrany ne
\begin{itemize}
\item tah může být uzavřený (končí, kde začal), nebo otevřený
\end{itemize}
\item eulerovský tah
\begin{itemize}
\item eulerovský tah obsahuje všechny hrany a vrcholy grafu
\begin{itemize}
\item někdy se eulerovský tah definuje bez nutnosti obsahovat všechny vrcholy
\end{itemize}
\item graf je eulerovský $\equiv$ existuje v něm uzavřený eulerovský tah
\end{itemize}
\item věta: graf $G$ je eulerovský $\iff G$ je souvislý a každý jeho vrchol má sudý stupeň
\end{itemize}
\item Stromy - definice a základní vlastnosti (existence listů, počet hran stromu)
\begin{itemize}
\item strom je souvislý graf bez kružnic (= acyklický)
\item les je acyklický graf
\item list je vrchol stupně 1
\begin{itemize}
\item strom o jednom vrcholu (\quotedblbase{}pařez\textquotedblleft{}) nemá žádný list
\end{itemize}
\item lemma: každý strom s aspoň 2 vrcholy má aspoň 1 list (respektive aspoň dva listy)
\item lemma: je-li $v$ list grafu $G$, pak $G$ je strom, právě když $G-v$ je strom
\item strom $G$ má $|V(G)|-1$ hran (Eulerova formule)
\end{itemize}
\item Ekvivalentní charakteristiky stromů
\begin{enumerate}
\item G je souvislý a acyklický (strom)
\item mezi vrcholy $u, v$ existuje právě jedna cesta (jednoznačně souvislý)
\item G je souvislý a po smazání libovolné jedné hrany už nebude souvislý (minimální souvislý)
\item G je acyklický a po přidání libovolné jedné hrany vznikne cyklus (maximální acyklický)
\item G je souvislý a platí pro něj Eulerova formule $|E(G)|=|V(G)|-1$
\end{enumerate}
\item Rovinné grafy - definice a základní pojmy (rovinný graf a rovinné nakreslení grafu, stěny)
\begin{itemize}
\item rovinné nakreslení grafu = nakreslení grafu, aby se hrany nekřížily
\begin{itemize}
\item vrcholy = body v rovině (navzájem různé)
\item hrany = křivky, které se neprotínají a jejich společnými body jsou jejich společné vrcholy
\end{itemize}
\item stěny nakreslení
\begin{itemize}
\item části, na které nakreslení grafu rozděluje rovinu
\item stěnou je i vnější stěna (zbytek roviny)
\item hranice stěny - skládá se z hran
\item hranice stěny je nakreslení uzavřeného sledu
\end{itemize}
\item graf je rovinný, pokud má alespoň jedno rovinné nakreslení
\item topologický graf je uspořádaná dvojice (graf, nakreslení)
\item $K_5$ a $K_{3,3}$ nejsou rovinné
\begin{itemize}
\item graf je rovinný, právě když žádný jeho podgraf není izomorfní s nějakým dělením grafu $K_5$ nebo $K_{3,3}$ (tzn. podgraf může mít podrozdělené hrany - \quotedblbase{}$K_5$\textquotedblleft{} s libovolně podrozdělenými hranami pořád nebude rovinná)
\end{itemize}
\end{itemize}
\item Eulerova formule a maximální počet hran rovinného grafu (důkaz a použití)
\begin{itemize}
\item věta
\begin{itemize}
\item nechť $G$ je souvislý graf nakreslený do roviny, $v:=|V(G)|, e:=|E(G)|, f:=$ počet stěn nakreslení
\item potom $v+f=e+2$ (Eulerova formule)
\end{itemize}
\item důkaz: indukcí podle $e$
\begin{itemize}
\item $e=v-1$ (G je strom)
\begin{itemize}
\item $f=1$
\item $v+1=v-1+2\quad\checkmark$
\end{itemize}
\item $e-1\to e$
\begin{itemize}
\item mějme graf $G$ s $e$ hranami
\item zvolím si libovolnou hranu $x$ na kružnici
\item $G':=G-x$
\item $v'=v,\quad e'=e-1,\quad f'=f-1$
\item z IP: $v'+f'=e'+2$
\item po dosazení: $v+f-1=e-1+2$
\item k oběma stranám přičteme jedničku
\item $v+f=e+2\quad\square$
\end{itemize}
\end{itemize}
\item věta: maximální rovinný graf je triangulace
\begin{itemize}
\item dokud hranice stěny není trojúhelník, je tam nějaká dvojice nespojených vrcholů
\end{itemize}
\item počet hran maximálního rovinného grafu
\begin{itemize}
\item každá stěna přispěje třemi hranami
\item každá hrana patří ke dvěma stěnám
\item počítáme \quotedblbase{}strany hran\textquotedblleft{}: $3f=2e$
\item $f={2\over 3}e$
\item $v+{2\over 3}e=e+2$
\item $e=3v-6$
\end{itemize}
\item věta: v každém rovinném grafu s aspoň 3 vrcholy je $|E|\leq 3|V|-6$
\item důkaz
\begin{itemize}
\item doplníme do $G$ hrany, až získáme maximální rovinný $G'$
\item $e'=3v-6$ (vrcholy nepřidáváme)
\item $e\leq e'=3v-6$
\end{itemize}
\item důsledek
\begin{itemize}
\item průměrný stupeň vrcholu v rovinném grafu je menší než 6
\begin{itemize}
\item $\sum \text{deg}(\xi)=2e\leq 6v-12$
\item průměrný stupeň $\leq {6v-12\over v}\lt 6$
\end{itemize}
\end{itemize}
\item počet hran a vrchol nízkého stupně v rovinných grafech bez trojúhelníků
\begin{itemize}
\item maximální rovinné grafy bez trojúhelníků mají stěny čtvercové, pětiúhelníkové, nebo to může být strom ve tvaru hvězdy
\item pro čtvercově stěny platí $4f=2e$, pro pětiúhelníkové $5f=2e$
\item obecně $4f\leq 2e\to f\leq \frac 12 e$
\item $v+\frac 12e\geq e+2$
\item $e\leq 2v-4$
\item průměrný stupeň $\leq {4v-8\over v} \lt$ 4
\item existuje vrchol stupně max. 3
\end{itemize}
\end{itemize}
\item Barevnost grafů - definice dobrého obarvení, vztah barevnosti a klikovosti grafu
\begin{itemize}
\item obarvení grafu $G$ $k$ barvami ($k$-obarvení) je $c:V(G)\to[k]$ t. ž. kdykoli $\lbrace x,y\rbrace\in E(G)$, pak $c(x)\neq c(y)$
\item barevnost $\chi(G)$ grafu $G:=\text{min }k:\exists$ k-obarvení grafu $G$
\item pozorování: kdykoli $H\subseteq G$, pak $\chi(H)\leq \chi(G)$
\item barevnost některých grafů
\begin{itemize}
\item úplné grafy \dots{} $\chi(K_n)=n$
\item cesty \dots{} $\chi(P_n)=2$ pro $n\geq 1$
\item sudé kružnice \dots{} $\chi(C_{2k})=2$
\item liché kružnice \dots{} $\chi(C_{2k+1})=3$
\item stromy jsou 2-obarvitelné
\item pro rovinné grafy existuje věta o čtyřech barvách (je to těsný horní odhad - existují rovinné grafy, které nelze obarvit třemi barvami, např. $K_4$)
\end{itemize}
\item věta: $\chi(G)\leq 2\iff G$ je bipartitní $\iff G$ neobsahuje lichou kružnici
\item definice: klikovost grafu $\kappa(G)$ je maximální $k$ takové, že v grafu jako podgraf existuje úplný graf $K_k$
\begin{itemize}
\item $\chi(G)\geq\kappa(G)$
\end{itemize}
\item klika v grafu je množina vrcholů taková, že každé dva jsou spojené hranou
\item nezávislá množina v grafu je množina vrcholů taková, že žádné dva vrcholy nejsou spojené hranou
\end{itemize}
\item Hranová a vrcholová souvislost grafů
\begin{itemize}
\item $F\subseteq E$ je hranový řez v $G$, pokud $G-F$ je nesouvislý
\begin{itemize}
\item kde $G-F\coloneqq(V,E\setminus F)$
\end{itemize}
\item graf je hranově $k$-souvislý, pokud neobsahuje žádný hranový řez velikosti menší než $k$
\item stupeň hranové souvislosti (nebo jen hranová souvislost) grafu $G$, značený $k_e(G)$, je největší $k$ takové, že $G$ je hranově $k$-souvislý
\begin{itemize}
\item pozorování: $k_e(G)=$ velikost nejmenšího hranového řezu v $G$
\end{itemize}
\item $A\subseteq V$ je vrcholový řez, pokud $G-A$ je nesouvislý
\begin{itemize}
\item kde $G-A=(V\setminus A,E\cap {V\setminus A\choose 2})$
\end{itemize}
\item graf je vrcholově $k$-souvislý, pokud má aspoň $k+1$ vrcholů a neobsahuje žádný vrcholový řez velikosti menší než $k$
\item vrcholová souvislost grafu $G$, značená $k_v(G)$, je největší $k$ takové, že $G$ je vrcholově $k$-souvislý
\begin{itemize}
\item pozorování: $k_v(K_n)=n-1$
\item pozorování: $G$ není úplný \dots{} $k_v(G)=$ velikost nejmenšího vrcholového řezu
\end{itemize}
\end{itemize}
\item Hranová a vrcholová verze Mengerovy věty
\begin{itemize}
\item Mengerova věta, hranová $xy$-verze
\begin{itemize}
\item pro dva různé vrcholy $x,y$ grafu $G$ platí, že $G$ obsahuje $k$ hranově disjunktních cest z $x$ do $y$ $\iff G$ neobsahuje hranový $xy$-řez velikosti menší než $k$ 
\end{itemize}
\item Mengerova věta, globální hranová verze
\begin{itemize}
\item graf $G$ je hranově $k$-souvislý $\iff$ mezi každými dvěma různými vrcholy existuje $k$ hranově disjunktních cest
\end{itemize}
\item definice: dvě cesty v $G$ z $x$ do $y$ jsou navzájem vnitřně vrcholově disjunktní (VVD), pokud nemají žádný společný vrchol kromě $x,y$
\item Mengerova věta, vrcholová $xy$-verze
\begin{itemize}
\item nechť $G=(V,E)$ je graf
\item nechť $x,y$ jsou různé \textbf{nesousední} vrcholy
\item nechť $k\in\mathbb N$
\item potom $G$ obsahuje $k$ navzájem VVD cest z $x$ do $y\iff G$ neobsahuje vrcholový $xy$-řez velikosti $\lt k$
\end{itemize}
\item Mengerova věta, globální vrcholová verze
\begin{itemize}
\item $G$ je vrcholově $k$-souvislý $\iff$ mezi každými dvěma různými vrcholy $x,y$ existuje $k$ navzájem VVD cest
\end{itemize}
\end{itemize}
\item Orientované grafy, silná a slabá souvislost
\begin{itemize}
\item orientovaný graf \dots{} $(V,E): E \subseteq V^2 \setminus \{(x,x)\mid x\in V\}$
\begin{itemize}
\item nepovolíme smyčky (v podstatě zakážeme diagonálu na relaci)
\end{itemize}
\item podkladový graf je neorientovaný graf založený na tom původním orientovaném
\begin{itemize}
\item pro orientovaný $G=(V,E)$ existuje podkladový $G^0=(V,E^0)$, kde $\lbrace u,v\rbrace \in E^0\equiv (u,v)\in E\lor (v,u)\in E$
\end{itemize}
\item vstupní a výstupní stupeň $\text{deg}^\text{in},\text{deg}^\text{out}$ (hrany vedoucí do vrcholu / z vrcholu)
\item vrchol je vyvážený $\equiv \text{deg}^\text{in}(v)=\text{deg}^\text{out}(v)$
\item graf je vyvážený $\equiv$ všechny vrcholy jsou vyvážené
\item součet vstupních stupňů = součet výstupních stupňů = počet hran
\item orientovaný graf je slabě souvislý $\equiv$ jeho podkladový graf je souvislý
\item o. graf je silně souvislý $\equiv$ existuje orientovaná cesta mezi každými dvěma vrcholy
\item silná souvislost $\implies$ slabá souvislost
\item věta: pro orientovaný graf $G$ platí: $G$ je vyvážený a slabě souvislý $\iff$ $G$ je eulerovský $\iff$ $G$ je vyvážený a silně souvislý
\end{itemize}
\item Toky v sítích: definice sítě a toku v ní
\begin{itemize}
\item tokovou síť tvoří
\begin{itemize}
\item $V$ \dots{} množina vrcholů
\item $E$ \dots{} množina orientovaných hran, $E\subseteq V\times V$
\item $z\in V$ \dots{} zdroj
\item $s\in V\setminus\set{z}$ \dots{} spotřebič / stok
\item $c:E\to [0,+\infty)$ \dots{} $c(e)$ je kapacita hrany $e$
\end{itemize}
\item definice umožňuje mít hrany oběma směry, připouští také smyčky, ty však z hlediska toků v sítích nejsou nijak užitečné
\item poznámka: ze stoku může vycházet orientovaná hrana, podobně do zdroje může vést hrana
\item tok v síti $(V,E,z,s,c)$ je funkce $f:E\to[0,+\infty)$ splňující
\begin{itemize}
\item $\forall e\in E:0\leq f(e)\leq c(e)$
\item $\forall x\in V\setminus \set{z,s}:$ součet toků do vrcholu $x$ se rovná součtu toků z vrcholu (formálně pomocí sum) \dots{} Kirchhoffův zákon
\end{itemize}
\end{itemize}
\item Existence maximálního toku (bez důkazu)
\begin{itemize}
\item velikost toku $f$ v síti $(V,E,z,s,c)$ je $w(f)\coloneqq f[\text{Out}(z)]-f[\text{In}(z)]$
\item kde $\mathrm{Out}(z)$ a $\mathrm{In}(z)$ jsou množiny hran do, respektive ze zdroje a $f[A]\coloneqq\sum_{e\in A}f(e)$ pro $A\subseteq E$
\item maximální tok je tok, který má největší velikost
\item fakt: v každé tokové síti existuje maximální tok
\item poznámka: obecně na přednášce uvažujeme konečné tokové sítě, grafy apod.
\item idea důkazu
\begin{itemize}
\item mějme síť, očíslujme hrany (od $1$ do $m$)
\item tok $f$ reprezentujme jako uspořádanou $m$-tici hodnot
\item množina všech toků je uzavřená a omezená podmnožina $\mathbb R^m$, je tedy kompaktní
\item funkce, která toku $f$ přiřadí $w(f)$, je spojitá
\item víme z analýzy, že spojitá funkce na kompaktní množině nabývá maxima
\end{itemize}
\item řez v síti $(V,E,z,s,c)$ je množina hran $R\subseteq E$ taková, že každá orientovaná cesta ze $z$ do $s$ obsahuje aspoň jednu hranu $R$
\item Minimaxová věta o toku a řezu: nechť $f_\text{max}$ je maximální tok a $R_\text{min}$ je minimální řez v $(V,E,z,s,c)$, potom $w(f_\text{max})=c(R_\text{min})$
\end{itemize}
\item Princip hledání maximálního toku v síti s~celočíselnými kapacitami (například pomocí Ford-Fulkersonova algoritmu)
\begin{itemize}
\item rezerva hrany $uv$
\begin{itemize}
\item $r(uv)\coloneqq c(uv)-f(uv)+f(vu)$
\end{itemize}
\item hrana je nasycená $\equiv$ má nulovou rezervu
\item hrana je nenasycená $\equiv$ má kladnou rezervu
\item cesta je nenasycená $\equiv$ žádná její hrana není nasycená / všechny mají kladné rezervy
\item algoritmus
\begin{itemize}
\item iterujeme, dokud existuje nenasycená cesta ze zdroje do stoku
\item spočítáme rezervu celé cesty (minimum přes rezervy hran cesty)
\item pro každou hranu upravíme tok - v~protisměru odečteme co nejvíc, zbytek přičteme po směru
\end{itemize}
\item pro celočíselné kapacity vrátí celočíselný tok
\item racionální kapacity převedeme na celočíselné
\item pro iracionální kapacity se může rozbít
\end{itemize}
\end{itemize}
\subsection{5. Pravděpodobnost a statistika}
\begin{itemize}
\item Pravděpodobnostní prostor, náhodné jevy, pravděpodobnost (definice, příklady)
\begin{itemize}
\item $\Omega$ \dots{} množina elementárních jevů
\item $\mathcal F\subseteq\mathcal P(\Omega)$ je prostor jevů, pokud\dots{}
\begin{itemize}
\item $\emptyset,\Omega\in\mathcal F$
\item $A\in\mathcal F\implies A^c=\Omega\setminus A\in\mathcal F$
\item $A_1,A_2,\ldots\in\mathcal F\implies\bigcup A_i\in \mathcal F$
\end{itemize}
\item $P:\mathcal F\to[0,1]$ je pravděpodobnost, pokud\dots{}
\begin{itemize}
\item $P(\Omega)=1$
\item $P(\bigcup A_i)=\sum P(A_i)$ pro $A_1,A_2,\ldots\in \mathcal F$ po dvou disjunktní
\end{itemize}
\item příklady pravděpodobnostních prostorů
\begin{itemize}
\item klasický - tedy konečný s uniformní pravděpodobností, např. $\Omega=[6]$ (hod kostkou)
\begin{itemize}
\item příklad elementárního jevu \dots{} padla šestka
\item příklad jevu \dots{} padlo sudé číslo
\end{itemize}
\item diskrétní - dovolíme cinknutou kostku (nevyžadujeme uniformní pravděpodobnost)
\item geometrický - např. \quotedblbase{}náhodně\textquotedblleft{} střílíme na terč, takže šance, že zasáhneme konkrétní oblast, je úměrná obsahu oblasti (ale může být i ve více než dvou dimenzích)
\item spojitý - schopný střelec se trefuje spíše do středu terče (opět dovolíme i jiné než uniformní rozdělení)
\end{itemize}
\end{itemize}
\item Základní pravidla pro počítání s pravděpodobností
\begin{itemize}
\item $P(A)+P(A^c)=1$
\item $A\subseteq B\implies P(B\setminus A)=P(B)-P(A)\implies P(A)\leq P(B)$
\item $P(A\cup B)=P(A)+P(B)-P(A\cap B)$
\item $P(A_1\cup A_2\cup\dots)\leq\sum_i P(A_i)$ \dots{} subaditivita, Booleova nerovnost
\end{itemize}
\item Nezávislost náhodných jevů, podmíněná pravděpodobnost
\begin{itemize}
\item jevy $A,B\in\mathcal F$ jsou nezávislé, pokud $P(A\cap B)=P(A)\cdot P(B)$
\item můžeme uvažovat i větší množiny jevů
\begin{itemize}
\item jevy v množině jsou (vzájemně) nezávislé, pokud pro každou konečnou podmnožinu platí, že $P(\bigcap)=\prod P$
\item pokud podmínka platí jen pro dvouprvkové podmnožiny, jsou jevy \textit{po dvou nezávislé}
\end{itemize}
\item pokud $A,B\in\mathcal F$ a $P(B)\gt 0$, definujeme podmíněnou pravděpodobnost $A$ při $B$ jako $P(A\mid B)=\frac{P(A\cap B)}{P(B)}$
\item zjevně $P(A\cap B)=P(A)\cdot P(B\mid A)$
\item $P(A_1\cap A_2\cap \dots\cap A_n)=P(A_1)P(A_2\mid A_1)P(A_3\mid A_1\cap A_2)\dots P(A_n\mid A_1\cap\dots\cap A_{n-1})$
\begin{itemize}
\item lze ukázat indukcí (nebo neformálně rozepsáním členů vpravo a vykrácením)
\end{itemize}
\item věta o úplné pravděpodobnosti
\begin{itemize}
\item mějme $B_1,B_2,\dots$ rozklad $\Omega$
\item $P(A)=\sum_iP(B_i)\cdot P(A\mid B_i)$
\begin{itemize}
\item sčítance s $P(B_i)=0$ považujeme za nulové
\end{itemize}
\end{itemize}
\end{itemize}
\item Bayesův vzorec
\begin{itemize}
\item mějme $B_1,B_2,\dots$ rozklad $\Omega$
\item $P(B_j\mid A)=\frac{P(B_j)\cdot P(A\mid B_j)}{\sum_i P(B_i)\cdot P(A\mid B_i)}$
\end{itemize}
\item Náhodné veličiny a jejich rozdělení
\begin{itemize}
\item náhodná veličina je funkce, která přiřazuje každému elementárnímu náhodnému jevu nějakou (obvykle číselnou) hodnotu
\item pro diskrétní náhodnou veličinu $X$ definujeme pravděpodobnostní funkci $p_X(x)=P(\set{X=x})$
\item spojité náhodné veličiny obvykle popisujeme pomocí distribuční funkce a hustoty
\begin{itemize}
\item distribuční funkce náhodné veličiny $X$ je funkce $F_X:\mathbb R\to [0,1]$ definovaná předpisem $F_X(x):=P(X\leq x)$
\item zjevně $P(a\lt X\leq b)=F_X(b)-F_X(a)$
\item vlastnosti
\begin{itemize}
\item $F_X$ je neklesající
\item $\lim_{x\to+\infty} F_X(x)=1$
\item $\lim_{x\to-\infty} F_X(x)=0$
\item $F_X$ je zprava spojitá
\end{itemize}
\item náhodná veličina $X$ se nazývá spojitá, pokud existuje nezáporná reálná funkce $f_X$ tak, že $F_X(x)=P(X\leq x)=\int_{-\infty}^x f_X(t)\text dt$
\item hustota $f_X$ \dots{} \quotedblbase{}limita histogramů\textquotedblleft{}
\item zjevně $\int_{-\infty}^\infty f_X=1$
\end{itemize}
\end{itemize}
\item Střední hodnota - linearita střední hodnoty, střední hodnota součinu nezávislých veličin, Markovova nerovnost
\begin{itemize}
\item $\mathbb E(X)=\sum_{x\in\text{Im}(X)}x\cdot P(X=x)$
\begin{itemize}
\item nebo $\mathbb E(X)=\int_{-\infty}^{\infty} x f_X(x)\text dx$
\end{itemize}
\item pozorování: $\mathbb E(X)=\sum_{\omega\in\Omega}X(\omega)\cdot P(\set{\omega})$
\item pravidlo naivního statistika
\begin{itemize}
\item $\mathbb E(g(X))=\sum_{x\in\text{Im}(X)}g(x)P(X=x)$
\item nebo $\mathbb E(g(X))=\int_{-\infty}^\infty g(x)f_X(x)\text dx$
\end{itemize}
\item linearita $\mathbb E(aX+b)=a\mathbb E(X)+b$ plyne z PNS pro funkci $ax+b$
\item podmíněná střední hodnota $\mathbb E(X\mid B)=\sum_{x\in\text{Im}(X)}x\cdot P(X=x\mid B)$
\item věta o celkové střední hodnotě: $\mathbb E(X)=\sum_i P(B_i)\cdot \mathbb E(X\mid B_i)$
\item $\mathbb E(g(X,Y))=\sum_x\sum_y g(x,y) P(X=x\land Y=y)$
\item $\mathbb E(aX+bY)=a\mathbb E(X)+b\mathbb E(Y)$
\item pro nezávislé $X,Y$ platí $\mathbb E(XY)=\mathbb E(X)\mathbb E(Y)$
\item Markovova nerovnost
\begin{itemize}
\item pro $X\geq 0$ a $a\gt0$ platí $P(X\geq a)\leq\frac{\mathbb E(X)}a$
\end{itemize}
\end{itemize}
\item Rozptyl - definice, vzorec pro rozptyl součtu (závislých či nezávislých veličin)
\begin{itemize}
\item rozptyl \dots{} $\text{var}(X)=\mathbb E((X-\mathbb EX)^2)$
\item směrodatná odchylka \dots{} $\sigma_X=\sqrt{\text{var}(X)}$
\item věta: $\text{var}(X)=\mathbb E(X^2)-\mathbb E(X)^2$
\item počítání s rozptylem
\begin{itemize}
\item $\mathrm{var}(aX+b)=a^2\mathrm{var}(X)$
\end{itemize}
\item rozptyl součtu nezávislých veličin je součet rozptylů
\begin{itemize}
\item $\mathrm{var}(X+Y)=\mathrm{var}(X)+\mathrm{var}(Y)$
\end{itemize}
\item kovariance a její vlastnosti
\begin{itemize}
\item definice: $\text{cov}(X,Y)=\mathbb E((X-\mathbb EX)(Y-\mathbb EY))$
\item věta: $\text{cov}(X,Y)=\mathbb E(XY)-\mathbb E(X)\mathbb E(Y)$
\item pro nezávislé $X,Y$ platí $\text{cov}(X,Y)=0$
\item zjevně $\text{cov}(X,X)=\text{var}(X)$
\item $\text{cov}(aX+bY,Z)=a\text{cov}(X,Z)+b\text{cov}(Y,Z)$
\item korelace
\begin{itemize}
\item $\rho(X,Y)=\frac{\text{cov}(X,Y)}{\sigma_X\cdot\sigma_Y}$
\end{itemize}
\end{itemize}
\item rozptyl součtu náhodných veličin (závislých i nezávislých)
\begin{itemize}
\item nechť $X=\sum_{i=1}^n X_i$
\item pak $\text{var}(X)=\sum_{i=1}^n\sum_{j=1}^n\text{cov}(X_i,X_j)$
\end{itemize}
\end{itemize}
\item Práce s konkrétními rozděleními: geometrické, binomické, Poissonovo, normální, exponenciální
\begin{itemize}
\item Bernoulliho/alternativní rozdělení $\text{Ber}(p)$
\begin{itemize}
\item počet úspěchů při jednom pokusu (kde $p$ je pravděpodobnost úspěchu)
\item $p_X(1)=p$
\item $p_X(0)=1-p$
\item jinak $p_X(x)=0$
\item $\mathbb E(X)=p$
\item $\text{var}(X)=p(1-p)$
\end{itemize}
\item geometrické rozdělení $\text{Geom}(p)$
\begin{itemize}
\item při kolikátém pokusu poprvé uspějeme
\item $p_X(k)=(1-p)^{k-1}\cdot p$
\item $\mathbb E(X)=1/p$
\item $\text{var}(X)=\frac{1-p}{p^2}$
\end{itemize}
\item binomické rozdělení $\text{Bin}(n,p)$
\begin{itemize}
\item počet úspěchů při $n$ nezávislých pokusech
\item $p_X(k)={n\choose k}p^k(1-p)^{n-k}$
\item $\mathbb E(X)=np$
\item $\text{var}(X)=np(1-p)$
\end{itemize}
\item Poissonovo rozdělení $\text{Pois}(\lambda)$
\begin{itemize}
\item počet doručených zpráv za časový úsek, $\lambda$ je průměrná hodnota
\item $p_X(k)=\frac{\lambda^k}{k!}e^{-\lambda}$
\item $\text{Pois}(\lambda)$ je limitou $\text{Bin}(n,\lambda/n)$ pro $n\to\infty$
\item $\mathbb E(X)=\lambda$
\item $\text{var}(X)=\lambda$
\end{itemize}
\item uniformní $U(a,b)$
\begin{itemize}
\item $f_X(x)=\frac 1{b-a}$
\item $F_X(x)=\frac{x-a}{b-a}$
\item $\mathbb E(X)=\frac{a+b}2$
\item $\text{var}(X)=\frac{(b-a)^2}{12}$
\end{itemize}
\item exponenciální $\text{Exp}(\lambda)$
\begin{itemize}
\item $f_X(x)=\lambda e^{-\lambda x}$ pro $x\geq 0$
\item $F_X(x)=1-e^{-\lambda x}$ pro $x\geq 0$
\item $\mathbb E(X)=1/\lambda$
\item $\text{var}(X)=1/\lambda^2$
\end{itemize}
\item standardní normální rozdělení
\begin{itemize}
\item $N(0,1)$
\item $f_X(x)=\varphi(x)=\frac1{\sqrt{2\pi}} e^{-x^2/2}$
\end{itemize}
\item obecné normální rozdělení
\begin{itemize}
\item $N(\mu,\sigma^2)$
\item $f_X(x)=\frac 1{\sigma\sqrt{2\pi}}e^{-\frac12(\frac{x-\mu}{\sigma})^2}$
\item máme-li $X\sim N(\mu,\sigma^2)$, pak pro $Z=\frac{X-\mu}\sigma$ platí $Z\sim N(0,1)$
\item pro normální nezávislé náhodné veličiny $X_i\sim N(\mu_i,\sigma_i^2)$ (nechť je jich konečně) je i součet normální náhodná veličina $\sum X_i\sim N(\sum\mu_i,\sum\sigma_i^2)$
\item pravidlo $3\sigma$
\begin{itemize}
\item $P(\mu-\sigma\lt X\lt\mu+\sigma)\doteq 0.68$
\item $P(\mu-2\sigma\lt X\lt\mu+2\sigma)\doteq 0.95$
\item $P(\mu-3\sigma\lt X\lt\mu+3\sigma)\doteq 0.997$
\end{itemize}
\end{itemize}
\end{itemize}
\item Limitní věty: zákon velkých čísel
\begin{itemize}
\item mějme $X_1,X_2,\dots$ stejně rozdělené nezávislé náhodné veličiny
\item $\bar X_n=(X_1+\dots+X_n)/n$ \dots{} výběrový průměr
\item silný zákon velkých čísel
\begin{itemize}
\item $\lim_{n\to\infty}\bar X_n=\mu$ skoro jistě (s pravděpodobností 1)
\item použití: Monte Carlo integrování kruhu
\end{itemize}
\item slabý zákon velkých čísel (zlepšení přesnosti opakovaným měřením)
\begin{itemize}
\item $\forall\varepsilon\gt 0:\lim_{n\to\infty} P(|\bar X_n-\mu|\gt\varepsilon)=0$
\item říkáme, že $\bar X_n$ konverguje k $\mu$ v pravděpodobnosti, píšeme $\bar X_n\xrightarrow P\mu$
\end{itemize}
\end{itemize}
\item Limitní věty: centrální limitní věta
\begin{itemize}
\item nechť $X_1,X_2,\dots$ jsou stejně rozdělené se střední hodnotou $\mu$ a rozptylem $\sigma^2$
\item označme $Y_n=\frac{(X_1+\dots+X_n)-n\mu}{\sigma\sqrt{n}}$
\item pak $Y_n\xrightarrow d N(0,1)$
\begin{itemize}
\item $Y_n$ konverguje v distribuci k $N(0,1)$
\item tzn. $\lim_{n\to\infty} F_{Y_n}(x)=\Phi(x)$
\end{itemize}
\item CLV se hodí k aproximaci distribuce součtu nebo průměru velkého počtu náhodných veličin normálním rozdělením
\begin{itemize}
\item takže můžeme provádět bodové a intervalové odhady i tam, kde data nejsou normálně rozdělená, ale známe rozptyl
\end{itemize}
\end{itemize}
\item Bodové odhady - alespoň jedna metoda pro jejich tvorbu, vlastnosti
\begin{itemize}
\item poznámka: neodhadujeme přímo $\theta$, nýbrž $g(\theta)$, aby to bylo obecnější
\item pro náhodný výběr $X_1,\dots,X_n\sim F_\theta$ a libovolnou funkci $g$ nazveme bodový odhad $\hat\theta_n$\dots{}
\begin{itemize}
\item nevychýlený/nestranný, pokud $\mathbb E(\hat\theta_n)=g(\theta)$
\item asymptoticky nevychýlený, pokud $\lim_{n\to\infty}\mathbb E(\hat\theta_n)=g(\theta)$
\item konzistentní, pokud $\hat\theta_n\xrightarrow P g(\theta)$
\end{itemize}
\item dále definujeme
\begin{itemize}
\item vychýlení \dots{} $\text{bias}(\hat\theta_n)=\mathbb E(\hat\theta_n)-\theta$
\item střední kvadratickou chybu \dots{} $\text{MSE}(\hat\theta_n)=\mathbb E((\hat\theta_n-\theta)^2)$
\end{itemize}
\item věta: $\text{MSE}(\hat\theta_n)=\text{bias}(\hat\theta_n)^2+\text{var}(\hat\theta_n)$
\item metoda momentů
\begin{itemize}
\item $r$-tý moment $X$ \dots{} $\mathbb E(X^r)=m_r(\theta)$
\item $r$-tý výběrový moment \dots{} $\frac1n\sum_{i=1}^n X_i^r=\widehat{m_r}$
\begin{itemize}
\item konzistentní nevychýlený odhad pro $r$-tý moment
\end{itemize}
\item nalezneme $\theta$ takové, že $m_r(\theta)=\widehat{m_r}$
\item typicky stačí použít první moment, dostaneme nějakou rovnici
\item $m_1=\mu$
\item $m_2=\mathbb E(X^2)=\text{var}(X)+(\mathbb EX)^2=\sigma^2+\mu^2$
\end{itemize}
\item metoda maximální věrohodnosti
\begin{itemize}
\item $\hat\theta_{ML}=\text{argmax}_\theta\;p(x;\theta)$
\begin{itemize}
\item $\text{argmax}_\theta\;f(x;\theta)$
\item abych se nemusel rozhodovat mezi $p$ a $f$, budu používat $L$
\end{itemize}
\item výpočetně jednodušší bude používat logaritmus $L$, který označíme $\ell$
\item příklad
\begin{itemize}
\item ve vzorku $k$ leváků z $n$ lidí, hledáme pravděpodobnost $\theta$, že je někdo levák
\item $L(x;\theta)=\theta^k(1-\theta)^{n-k}$
\item $\ell(x;\theta)=k\log\theta+(n-k)\log(1-\theta)$
\item $\ell'(x;\theta)=\frac k\theta-\frac{n-k}{1-\theta}$
\begin{itemize}
\item hledáme maximum, položíme derivaci rovnou nule (a zkontrolujeme krajní hodnoty)
\end{itemize}
\end{itemize}
\item podobně pro spojitý případ - \quotedblbase{}maximalizujeme\textquotedblleft{} rovnici pro pravděpodobnost konkrétního výběru
\end{itemize}
\end{itemize}
\item Intervalové odhady: metoda založená na aproximaci normálním rozdělením
\begin{itemize}
\item statistiky $D\leq H$ určují konfidenční interval o spolehlivosti $1-\alpha$, pokud $P(D\leq\theta\leq H)=1-\alpha$
\begin{itemize}
\item zkráceně $(1-\alpha)$-CI
\end{itemize}
\item interval budeme uvažovat ve tvaru $[x-\delta,x+\delta]$
\item postup
\begin{itemize}
\item máme nestranný bodový odhad $\hat\theta$ pro parametr $\theta$
\item $\hat\theta$ má normální rozdělení
\item $\hat\theta\pm z_{\alpha/2}\cdot\text{se}$ je $(1-\alpha)$-CI
\begin{itemize}
\item $z_{\alpha/2}:=\Phi^{-1}(1-\alpha/2)$
\item $\text{se}:=\sigma(\hat\theta)$
\end{itemize}
\end{itemize}
\end{itemize}
\item Testování hypotéz - základní přístup, chyby 1. a 2. druhu, hladina významnosti
\begin{itemize}
\item zvolíme testovou statistiku $T$
\item nulová hypotéza \dots{} defaultní, konzervativní model
\item alternativní hypotéza \dots{} alternativní model, \quotedblbase{}zajímavost\textquotedblleft{}
\item nulovou hypotézu buď zamítneme, nebo nezamítneme
\item chyba 1. druhu - chybné zamítnutí, \quotedblbase{}trapas\textquotedblleft{}
\item chyba 2. druhu - chybné přijetí, \quotedblbase{}promarněná příležitost\textquotedblleft{}
\item hladina významnosti $\alpha$ \dots{} pravděpodobnost chyby 1. druhu
\begin{itemize}
\item typicky se volí $\alpha=0.05$
\end{itemize}
\item $\beta$ \dots{} pravděpodobnost chyby 2. druhu
\item kritický obor \dots{} je to množina, kterou určíme před provedením testu; pokud se výsledek našeho testu bude nacházet v~kritickém oboru, zamítneme nulovou hypotézu
\begin{itemize}
\item tedy $\alpha=P(h(X)\in W; H_0)$
\end{itemize}
\item síla testu \dots{} $1-\beta$
\begin{itemize}
\item chceme co největší
\end{itemize}
\item $p$-hodnota \dots{} nejmenší $\alpha$ taková, že na hladině $\alpha$ zamítáme $H_0$
\begin{itemize}
\item takže zamítáme, pokud $p$-hodnota vyjde menší rovna naší zvolené $\alpha$
\item $p$-hodnota je vlastně pravděpodobnost, že při platnosti $H_0$ nabývá zvolená testová statistika $T$ naměřené hodnoty nebo hodnot ještě extrémnějších
\end{itemize}
\end{itemize}
\end{itemize}
\subsection{6. Logika}
\begin{itemize}
\item Znalost a práce se základními pojmy syntaxe výrokové a predikátové logiky (jazyk, otevřená a uzavřená formule, instance formule, apod.)
\begin{itemize}
\item výroková logika
\begin{itemize}
\item jazyk je určený neprázdnou množinou výrokových proměnných $\mathbb P$ (také jim říkáme prvovýroky nebo atomické výroky)
\item množina $\mathbb P$ může být konečná nebo i nekonečná (ale obvykle bude spočetná) a bude mít dané uspořádání
\end{itemize}
\item predikátová logika
\begin{itemize}
\item signatura je dvojice $\braket{\mathcal {R,F}}$, kde $\mathcal{R,F}$ jsou disjunktní množiny symbolů (relační a funkční, ty zahrnují konstantní) spolu s danými aritami a neobsahují symbol $=$
\begin{itemize}
\item signaturu ale často uvádíme jako prostý výčet symbolů
\end{itemize}
\item do jazyka patří\dots{}
\begin{itemize}
\item spočetně mnoho proměnných
\item relační, funkční a konstantní symboly ze signatury a symbol $=$, jde-li o jazyk s rovností
\item univerzální a existenční kvantifikátory pro každou proměnnou
\item symboly pro logické spojky a závorky
\end{itemize}
\item např. jazyk uspořádání zapíšeme jako \quotedblbase{}$L=\braket{\leq}$ s rovností\textquotedblleft{}
\begin{itemize}
\item tzn. napíšeme signaturu a také uvedeme, zda jde o jazyk s rovností
\end{itemize}
\item struktura v signatuře $\braket{\mathcal{R,F}}$ je trojice $\mathcal A=\braket{A,\mathcal{R^A,F^A}}$, kde\dots{}
\begin{itemize}
\item $A$ je neprázdná množina, říkáme jí doména (také universum)
\item $\mathcal {R^A}$ je množina interpretací jednotlivých relačních symbolů (kde interpretace $n$-árního relačního symbolu je množina uspořádaných $n$-tic prvků z $A$)
\item $\mathcal {F^A}$ je množina interpretací jednotlivých funkčních symbolů (kde intepretace $n$-árního funkčního symbolu odpovídá funkci přiřazující $n$-tici prvků z $A$ jeden prvek z $A$)
\begin{itemize}
\item speciálně pro konstantní symbol $c\in\mathcal F$ je $c^\mathcal A\in A$
\end{itemize}
\end{itemize}
\item termy jazyka $L$ jsou konečné nápisy definované induktivně
\begin{itemize}
\item proměnná je term
\item konstantní symbol je term
\item pro $n$-ární funkční symbol $f$ a termy $t_1,\dots,t_n$ je nápis $f(t_1,\dots,t_n)$ také term
\end{itemize}
\item atomická formule jazyka $L$ je nápis $R(t_1,\dots,t_n)$, kde $R$ je $n$-ární relační symbol z $L$ (v jazyce s rovností to může být i rovnost) a $t_i$ je $L$-term
\item formule jazyka $L$ jsou konečné nápisy definované induktivně
\begin{itemize}
\item atomická formule je formule
\item negace formule je formule
\item spojením dvou formulí binární logickou spojkou dostaneme formuli
\item přidáním kvantifikátoru před formuli dostaneme formuli
\end{itemize}
\item výskyt proměnné je vázaný, pokud jí odpovídá nějaký kvantifikátor, jinak je výskyt volný
\item proměnná je volná, pokud má volný výskyt, je vázaná, pokud má vázaný výskyt
\item formule je otevřená, neobsahuje-li žádný kvantifikátor
\item formule je uzavřená (= sentence), pokud nemá žádnou volnou proměnnou
\item term $t$ je substituovatelný za proměnnou $x$ ve formuli $\varphi$, pokud po simultánním nahrazení všech volných výskytů $x$ ve $\varphi$ za $t$ nevznikne ve $\varphi$ žádný vázaný výskyt proměnné z $t$
\begin{itemize}
\item v tom případě říkáme vzniklé formuli instance $\varphi$ vzniklá substitucí $t$ za $x$ a označujeme ji $\varphi(x/t)$
\end{itemize}
\end{itemize}
\end{itemize}
\item Prenexní tvary formulí predikátové logiky
\begin{itemize}
\item PNF = kvantifikátory jsou před formulí, formule se dělí na kvantifikátorový prefix a otevřené jádro
\item převod na PNF - postupně kvantifikátory vytahuju (podle pravidel), v případě potřeby přejmenuju proměnnou, tím vznikne varianta (pokud je v druhé části formule volná proměnná se stejným názvem)
\item pravidlo převodu: pokud vytahujeme kvantifikátor z negace, musíme ho \quotedblbase{}přepnout\textquotedblleft{} (pozor: u implikace je levá strana jakoby negovaná)
\end{itemize}
\item Znalost základních normálních tvarů (CNF, DNF, PNF)
\begin{itemize}
\item PNF = kvantifikátory jsou před formulí, formule se dělí na kvantifikátorový prefix a otevřené jádro
\item tvary výroků
\begin{itemize}
\item literál je prvovýrok nebo negace prvovýroku
\begin{itemize}
\item pro literál $\ell$ označuje $\overline\ell$ literál k němu opačný (tedy jeho negaci)
\end{itemize}
\item klauzule je disjunkce literálů
\begin{itemize}
\item jednotková klauzule je samotný literál
\item prázdná klauzule je $\bot$
\end{itemize}
\item výrok je v konjunktivní normální formě (CNF), pokud je konjunkcí klauzulí
\begin{itemize}
\item prázdný výrok v CNF je $\top$
\end{itemize}
\item elementární konjunkce je konjunkce literálů
\begin{itemize}
\item jednotková elementární konjunkce je samotný literál
\item prázdná elementární konjunkce je $\top$
\end{itemize}
\item výrok je v disjunktivní normální formě (DNF), pokud je disjunkcí elementárních konjunkcí
\begin{itemize}
\item prázdný výrok v DNF je $\bot$
\end{itemize}
\item výrok je hornovský (v Hornově tvaru), pokud je konjunkcí hornovských klauzulí, tj. klauzulí obsahujících nejvýše jeden pozitivní literál
\end{itemize}
\item množinová reprezentace
\begin{itemize}
\item klauzule je konečná množina literálů
\begin{itemize}
\item prázdnou klauzuli označíme $\square$, není nikdy splněna
\end{itemize}
\item CNF formule je (konečná nebo nekonečná) množina klauzulí
\begin{itemize}
\item prázdná formule $\emptyset$ je vždy splněna
\end{itemize}
\end{itemize}
\end{itemize}
\item Převody na normální tvary
\begin{itemize}
\item sémantický převod
\begin{itemize}
\item najdeme modely výroku
\item pro DNF budeme uvažovat disjunkci elementárních konjunkcí, přičemž každá elementární konjunkce popíše jeden model
\item pro CNF musíme najít nemodely formule, každá klauzule zakáže jeden nemodel
\begin{itemize}
\item pokud je první nemodel $(0,0,0)$, pak první klauzule bude $(p\lor q\lor r)$
\end{itemize}
\end{itemize}
\item syntaktický převod
\begin{itemize}
\item přepíšeme implikaci a ekvivalenci pomocí konjunkce, disjunkce a negace
\item přesuneme negace dovnitř (k literálům)
\item použijeme distributivitu konjunkce a disjunkce, abychom jednu přesunuli dovnitř (podle toho, jestli chceme CNF nebo DNF)
\end{itemize}
\end{itemize}
\item Použití normálních tvarů pro algoritmy (SAT, rezoluce)
\begin{itemize}
\item problém Horn-SAT
\begin{itemize}
\item výrok je hornovský, pokud je konjunkcí hornovských klauzulí, tj. klauzulí obsahujících nejvýše jeden pozitivní literál
\item algoritmus
\begin{itemize}
\item pokud $\varphi$ obsahuje dvojici opačných jednotkových klauzulí, není splnitelný
\item pokud $\varphi$ neobsahuje žádnou jednotkovou klauzuli, je splnitelný, ohodnotíme všechny zbývající proměnné nulou
\item pokud $\varphi$ obsahuje jednotkovou klauzuli $\ell$, ohodnotíme literál $\ell$ hodnotou 1, provedeme jednotkovou propagaci a postup opakujeme
\end{itemize}
\item jednotková propagace pro $\ell=1$
\begin{itemize}
\item každou klauzuli obsahující $\ell$ odstraníme (protože je takto splněna)
\item $\overline\ell$ odstraníme ze všech klauzulí, které ho obsahují (protože $\overline\ell$ nemůže zajistit splnění dané klauzule)
\end{itemize}
\item Horn-SAT lze řešit v lineárním čase
\end{itemize}
\item algoritmus DPLL pro řešení SAT
\begin{itemize}
\item literál $\ell$ má čistý výskyt ve $\varphi$, pokud se $\ell$ vyskytuje ve $\varphi$ a opačný literál $\overline\ell$ se ve $\varphi$ nevyskytuje
\begin{itemize}
\item takový literál můžu nastavit na 1
\item to neovlivní splnitelnost výroku, ale zmenší to množinu modelů, které jsem schopen nalézt
\end{itemize}
\item algoritmus
\begin{itemize}
\item dokud $\varphi$ obsahuje jednotkovou klauzuli $\ell$, ohodnoť $\ell=1$ a proveď jednotkovou propagaci
\item dokud existuje literál $\ell$, který má ve $\varphi$ čistý výskyt, ohodnoť $\ell=1$ a odstraň klauzule obsahující $\ell$
\item pokud $\varphi$ neobsahuje žádnou klauzuli, je splnitelný
\item pokud $\varphi$ obsahuje prázdnou klauzuli, není splnitelný
\item jinak zvol dosud neohodnocenou výrokovou proměnnou $p$ a zavolej algoritmus rekurzivně na $\varphi\land p$ a na $\varphi\land\neg p$
\end{itemize}
\item algoritmus běží v exponenciálním čase
\end{itemize}
\item rezoluce
\begin{itemize}
\item použitím rezolučního pravidla z klauzulí $\varphi\lor p$ a $\psi\lor\neg p$ dostaneme $\varphi\lor\psi$
\item takhle spolu můžeme postupně rezolvovat různé klauzule z teorie
\item pokud nám vyjde prázdná klauzule, je teorie sporná
\item k dokázání výroku $\varphi$ v teorii $T$ hledáme rezoluční zamítnutí $T\cup\neg\varphi$
\begin{itemize}
\item nebo přímo rezoluční důkaz výroku $\varphi$ (tzn. na konci rezoluce nám vyjde $\varphi$)
\end{itemize}
\end{itemize}
\end{itemize}
\item Pojem modelu teorie
\begin{itemize}
\item model (jazyka) ve výrokové logice je pravdivostní ohodnocení proměnných
\begin{itemize}
\item model $v$ je modelem teorie $T$, pokud každý axiom z $T$ v modelu $v$ platí, píšeme $v\models T$
\end{itemize}
\item v predikátové logice: model jazyka $L$ nebo také $L$-struktura je libovolná struktura v signatuře jazyka $L$
\begin{itemize}
\item model teorie $T$ je $L$-struktura, ve které platí všechny axiomy teorie $T$
\end{itemize}
\end{itemize}
\item Pravdivost, lživost, nezávislost formule vzhledem k teorii, splnitelnost, tautologie, důsledek
\begin{itemize}
\item pravdivý výrok platí v každém modelu (jazyka/teorie)
\begin{itemize}
\item je pravdivý v logice = platí v logice = je tautologie
\item je pravdivý v $T$ = platí v $T$ = je důsledek $T$
\end{itemize}
\item lživý (sporný) výrok neplatí v žádném modelu
\item nezávislý výrok platí v nějakém modelu a neplatí v~jiném modelu
\item splnitelný výrok platí v nějakém modelu (tedy není lživý, je pravdivý nebo nezávislý)
\end{itemize}
\item Analýza výrokových teorií nad konečně mnoha prvovýroky
\begin{itemize}
\item nechť $T$ je bezesporná teorie nad $\mathbb P$
\begin{itemize}
\item $|\mathbb P|=n\in\mathbb N^+$ \dots{} počet proměnných v jazyce
\item $m=|M_{\mathbb P}(T)|$ \dots{} velikost množiny modelů teorie $T$
\end{itemize}
\item neekvivalentních výroků nad $\mathbb P$ je $2^{2^n}$
\begin{itemize}
\item každý výrok odpovídá jedné podmnožině množiny modelů (všech modelů je $2^n$)
\end{itemize}
\item neekvivalentních výroků nad $\mathbb P$ pravdivých/lživých v $T$ je $2^{2^n-m}$
\begin{itemize}
\item u pravdivých musí množina modelů obsahovat všech $m$ modelů, takže je zafixujeme (u lživých naopak nesmí obsahovat žádný z nich)
\end{itemize}
\item neekv. výroků nad $\mathbb P$ nezávislých v $T$ je $2^{2^n}-2\cdot 2^{2^n-m}$
\begin{itemize}
\item vezmeme všechny a odečteme pravdivé a lživé
\end{itemize}
\item neekv. jednoduchých extenzí teorie $T$ je $2^m$, z toho sporná je jedna
\begin{itemize}
\item podle toho, které modely ubereme
\end{itemize}
\item neekv. kompletních jednoduchých extenzí teorie $T$ je $m$
\begin{itemize}
\item každá odpovídá jednomu modelu
\end{itemize}
\item $T$-neekvivalentních výroků nad $\mathbb P$ je $2^m$
\item $T$-neekv. výroků nad $\mathbb P$ pravdivých/lživých (v $T$) je $1$
\item $T$-neekv. výroků nad $\mathbb P$ nezávislých (v $T$) je $2^m-2$
\end{itemize}
\item Extenze teorií - schopnost porovnat sílu teorií, konzervativnost, skolemizace
\begin{itemize}
\item extenze teorie $T$ je libovolná teorie $T'$ v jazyce $\mathbb P'\supseteq\mathbb P$ splňující $\text{Csq}_{\mathbb P}(T)\subseteq\text{Csq}_{\mathbb P'}(T')$
\begin{itemize}
\item kde $\text{Csq}_\mathbb P(T)$ je množina všech důsledků teorie (výroků/sentencí pravdivých v teorii $T$ v jazyce $\mathbb P$)
\end{itemize}
\item extenze je jednoduchá, pokud $\mathbb P'=\mathbb P$
\item extenze je konzervativní, pokud $\text{Csq}_{\mathbb P}(T)=\text{Csq}_{\mathbb P}(T')=\text{Csq}_{\mathbb P'}(T')\cap\text{VF}_\mathbb P$
\item sémantický význam (v řeči modelů)
\begin{itemize}
\item mějme teorii $T$ v jazyce $\mathbb P$ a teorii $T'$ v jazyce $\mathbb P'$, kde $\mathbb P\subseteq\mathbb P'$
\item $T'$ je jednoduchou extenzí $T$, právě když $\mathbb P'=\mathbb P$ a $M_\mathbb P(T')\subseteq M_\mathbb P(T)$
\item $T'$ je extenzí $T$, právě když $M_{\mathbb P'}(T')\subseteq M_{\mathbb P'}(T)$
\item $T'$ je konzervativní extenzí $T$, pokud je extenzí a navíc platí, že každý model $T$ lze nějak expandovat na model $T'$
\item $T'$ je extenzí $T$ a zároveň $T$ je extenzí $T'$, právě když $T'\sim T$ (jazyky a množiny modelů se rovnají)
\item kompletní jednoduché extenze $T$ jednoznačně až na ekvivalenci odpovídají modelům $T$
\end{itemize}
\item $T'$ je extenzí $T$ o definice, pokud vznikla z $T$ postupnou extenzí o definice relačních a funkčních (případně konstantních) symbolů
\item je-li $T'$ extenze teorie $T$ o definice, potom platí:
\begin{itemize}
\item každý model teorie $T$ lze jednoznačně expandovat na model $T'$
\item $T'$ je konzervativní extenze $T$
\item pro každou $L'$-formuli $\varphi'$ existuje $L$-formule $\varphi$ taková, že $T'\models\varphi'\leftrightarrow\varphi$
\end{itemize}
\item Skolemova varianta sentence vzniká z původní PNF sentence skolemizací
\begin{itemize}
\item skolemizace spočívá v tom, že tyto kroky iterujeme přes všechny existenční kvantifikátory $(\exists y_i)$
\begin{itemize}
\item odstraníme z prefixu existenční kvantifikátor $(\exists y_i$)
\item za proměnnou $y_i$ substituujeme term $f_i(x_1,\dots,x_{n_i})$, kde $x_1,\dots,x_{n_i}$ jsou proměnné, jejichž univerzální kvantifikátory předcházejí $(\exists y_i)$
\end{itemize}
\item začínáme s $L$-sentencí v PNF, jejíž všechny vázané proměnné jsou různé
\item dostaneme $L'$-sentenci v PNF, kde $L'$ je rozšíření $L$ o nové $n_i$-ární funkční symboly
\end{itemize}
\end{itemize}
\item Dokazatelnost - pojem formálního důkazu, zamítnutí; schopnost práce v některém z formálních dokazovacích systémů (např. tablo)
\begin{itemize}
\item důkaz, zamítnutí
\begin{itemize}
\item důkaz faktu, že v teorii $T$ platí výrok $\varphi$ je konečný syntaktický objekt vycházející z axiomů $T$ a výroku $\varphi$
\item důkaz konstruujeme syntakticky, aniž bychom se museli zabývat modely
\item po dokazovacím systému požadujeme dvě vlastnosti: korektnost (je-li výrok dokazatelný z $T$, je pravdivý v $T$) a úplnost (je-li výrok pravdivý v $T$, je dokazatelný z $T$)
\item důkaz říká, že $\varphi$ platí, zamítnutí říká, že $\varphi$ neplatí
\begin{itemize}
\item rezoluční důkaz $\varphi$ v teorii $S$ má v kořeni $\varphi$, rezoluční zamítnutí teorie $S$ má v kořeni $\square$
\end{itemize}
\end{itemize}
\item tablo důkaz
\begin{itemize}
\item konečné tablo z teorie $T$ je uspořádaný, položkami označkovaný strom zkonstruovaný aplikací konečně mnoha následujících pravidel:
\begin{itemize}
\item jednoprvkový strom označkovaný libovolnou položkou je tablo z teorie $T$
\item pro libovolnou položku $P$ na libovolné větvi $V$ můžeme na konec větve $V$ připojit atomické tablo pro položku $P$
\item na konec libovolné větve můžeme připojit položku $\text T\alpha$ pro libovolný axiom teorie $\alpha\in T$
\end{itemize}
\item tablo je konečné nebo nekonečné, každopádně vzniklo ve spočetně mnoha krocích
\item tablo pro položku $P$ je tablo s položkou $P$ v kořeni
\item tablo důkaz výroku $\varphi$ z teorie $T$ je sporné tablo z teorie $T$ s  položkou $\text F\varphi$ v~kořeni
\begin{itemize}
\item pokud existuje, je $\varphi$ tablo dokazatelný z $T$, píšeme $T\vdash\varphi$
\item podobně definujeme tablo zamítnutí s $\text T\varphi$ v kořeni, tablo zamítnutelnost se značí $T\vdash \neg\varphi$
\end{itemize}
\item tablo je sporné, pokud je každá jeho větev sporná
\item větev je sporná, pokud obsahuje položky $\text T\psi$ a $\text F\psi$ pro nějaký výrok $\psi$, jinak je bezesporná
\item tablo je dokončené, pokud je každá jeho větev dokončená
\item větev je dokončená\dots{}
\begin{itemize}
\item pokud je sporná
\item nebo pokud je každá její položka na této větvi redukovaná a pokud větev zároveň obsahuje položku s T pro každý axiom teorie
\end{itemize}
\item položka je redukovaná na dané větvi, pokud obsahuje pouze výrokovou proměnnou nebo se na dané větvi vyskytuje jako kořen atomického tabla (tedy došlo k jejímu rozvoji na dané větvi)
\item v predikátové logice navíc řešíme kvantifikátory - položky typu \textit{svědek} a \textit{všichni}
\end{itemize}
\end{itemize}
\item Věty o kompaktnosti a úplnosti výrokové a predikátové logiky - znění a porozumění významu; použití na příkladech, důsledky
\begin{itemize}
\item věta o kompaktnosti: teorie má model, právě když každá její konečná část má model
\item aplikace
\begin{itemize}
\item popíšeme požadovanou vlastnost nekonečného objektu pomocí (nekonečné) výrokové teorie
\item z konečné části teorie sestrojíme konečný podobjekt mající danou vlastnost
\end{itemize}
\item příklady
\begin{itemize}
\item spočetně nekonečný graf je bipartitní $\iff$ každý jeho konečný podgraf je bipartitní
\begin{itemize}
\item $T=\set{p_u\to\neg p_v\mid\set{u,v}\in E(G)}$
\end{itemize}
\item nestandardní model přirozených čísel
\begin{itemize}
\item vezmeme standardní model přirozených čísel (jeho teorii)
\item přidáme nový konstantní symbol
\item přidáme k teorii výrok, že je ta konstanta ostře větší než každý $n$-tý numerál
\item každá konečná část této nové teorie má model, takže i celá ta nová teorie má model
\end{itemize}
\end{itemize}
\item První věta o neúplnosti: Pro každou bezespornou rekurzivně axiomatizovanou extenzi $T$ Robinsonovy aritmetiky existuje sentence, která je pravdivá v $\underline{\mathbb N}$, ale není dokazatelná v $T$.
\begin{itemize}
\item Robinsonova aritmetika je teorie jazyka aritmetiky $L=\braket{S,+,\cdot,0,\leq}$ s rovností
\begin{itemize}
\item má osm (tradičně spíše sedm) axiomů, některé z nich:
\begin{itemize}
\item $S(x)\neq 0$
\item $S(x)=S(y)\to x=y$
\item $x+0=x$
\item $x+S(y)=S(x+y)$
\end{itemize}
\end{itemize}
\item $\underline{\mathbb N}=\braket{\mathbb N,S,+,\cdot,0,\leq}$ \dots{} standardní model přirozených čísel
\end{itemize}
\item důsledek 1.1: je-li $T$ rekurzivně axiomatizovaná extenze Robinsonovy aritmetiky a je-li navíc $\underline{\mathbb N}$ modelem teorie $T$, potom $T$ není kompletní
\begin{itemize}
\item pro sentenci, která není dokazatelná v $T$, nemůže být dokazatelná ani její negace (byl by to spor s její pravdivostí v $\underline{\mathbb N}$)
\end{itemize}
\item důsledek 1.2: teorie $\text{Th}(\underline{\mathbb N})$ není rekurzivně axiomatizovatelná
\begin{itemize}
\item kdyby byla, nemohla by být kompletní - ale ona kompletní je
\item značení: $\text{Th}(\underline{\mathbb N})$ \dots{} množina všech sentencí pravdivých v $\underline{\mathbb N}$ (tzv. teorie struktury $\underline{\mathbb N}$)
\end{itemize}
\item věta (Rosserův trik): v každé bezesporné rekurzivně axiomatizované extenzi Robinsonovy aritmetiky existuje nezávislá sentence - tedy taková není kompletní
\begin{itemize}
\item v podstatě plyne z důsledku 1.1, jen se zbavuje $\underline{\mathbb N}$ 
\end{itemize}
\item Druhá věta o neúplnosti: Pro každou bezespornou rekurzivně axiomatizovanou extenzi $T$ Peanovy aritmetiky platí, že $\mathit{Con_T}$ není dokazatelná v $T$.
\item $\mathit{Con_T}$ \dots{} sentence vyjadřující bezespornost (konzistence) teorie $T$
\begin{itemize}
\item $\mathit{Con_T}=\neg(\exists y)\mathit{Prf_{T}}(\underline{0=S(0)},y)$
\item $\mathit{Prf_{T}}(x,y)$ \dots{} $y$ je důkaz $x$ v $T$
\end{itemize}
\item důsledek 2.1: existuje model Peanovy aritmetiky (PA), ve kterém platí sentence $(\exists y)\mathit{Prf_{PA}}(\underline{0=S(0)},y)$
\item důsledek 2.2: existuje bezesporná rekurzivně axiomatizovaná extenze $T$ Peanovy aritmetiky, která dokazuje svou spornost, tj. taková, že $T\vdash\neg \mathit{Con_T}$
\item důsledek 2.3: je-li teorie množin ZFC bezesporná, nemůže být sentence $\mathit{Con_{ZFC}}$ v teorii ZFC dokazatelná
\end{itemize}
\item Rozhodnutelnost - pojem kompletnosti a její kritéria, význam pro rozhodnutelnost; příklady rozhodnutelných a nerozhodnutelných teorií
\begin{itemize}
\item teorie je sporná, jestliže\dots{} (ekvivalentně)
\begin{itemize}
\item v ní platí spor
\item nemá žádný model
\item v ní platí všechny výroky
\end{itemize}
\item teorie je bezesporná (splnitelná), pokud není sporná, tj. má nějaký model
\item teorie je kompletní jestliže není sporná a každý výrok (respektive sentence) je v ní pravdivý nebo lživý (tj. nemá žádné nezávislé výroky), ekvivalentně, pokud má právě jeden model (až na elementární ekvivalenci v predikátové logice)
\begin{itemize}
\item struktury $\mathcal{A,B}$ (v témž jazyce) jsou elementárně ekvivalentní, pokud v nich platí tytéž sentence (značíme $\mathcal A\equiv \mathcal B$)
\item zjevně $\mathcal A\equiv \mathcal B\iff\text{Th}(\mathcal A)=\text{Th}(\mathcal B)$
\end{itemize}
\item o teorii $T$ říkáme, že je\dots{}
\begin{itemize}
\item rozhodnutelná, pokud existuje algoritmus, který pro každou vstupní formuli $\varphi$ doběhne a odpoví, zda $T\models\varphi$
\item částečně rozhodnutelná, pokud existuje algoritmus, který pro každou vstupní formuli $\varphi$\dots{}
\begin{itemize}
\item pokud $T\vDash\varphi$, doběhne a odpoví \quotedblbase{}ano\textquotedblleft{}
\item pokud $T\nvDash\varphi$, buď nedoběhne, nebo doběhne a odpoví \quotedblbase{}ne\textquotedblleft{}
\end{itemize}
\end{itemize}
\item \quotedblbase{}rekurzivní\textquotedblleft{} pojmy
\begin{itemize}
\item teorie $T$ je rekurzivně axiomatizovaná, pokud existuje algoritmus, který pro každou vstupní formuli $\varphi$ doběhne a odpoví, zda $\varphi\in T$
\item třída $L$-struktur $K\subseteq M_L$ je rekurzivně axiomatizovatelná, pokud existuje rekurzivně axiomatizovaná $L$-teorie $T$ taková, že $K=M_L(T)$
\item teorie $T'$ je rekurzivně axiomatizovatelná, pokud je rekurzivně axiomatizovatelná třída jejích modelů, neboli pokud je $T'$ ekvivalentní nějaké rekurzivně axiomatizované teorii
\item řekneme, že teorie $T$ má rekurzivně spočetnou kompletaci, pokud (nějaká) množina (až na ekvivalenci) všech jednoduchých kompletních extenzí teorie $T$ je rekurzivně spočetná, tj. existuje algoritmus, který pro danou vstupní dvojici přirozených čísel $(i,j)$ vypíše na výstup $i$-tý axiom $j$-té extenze (v nějakém pevně daném uspořádání), nebo odpoví, že takový axiom už neexistuje
\end{itemize}
\item tvrzení: je-li $T$ rekurzivně axiomatizovaná, potom je částečně rozhodnutelná, je-li navíc kompletní, je rozhodnutelná
\item tvrzení: je-li $\mathcal A$ konečná struktura v konečném jazyce s rovností, potom je teorie $\text{Th}(\mathcal A)$ rekurzivně axiomatizovatelná
\item tvrzení: pokud je teorie $T$ rekurzivně axiomatizovaná a má rekurzivně spočetnou kompletaci, potom je $T$ rozhodnutelná
\item nerozhodnutelnost
\begin{itemize}
\item mějme formuli $\varphi$ ve tvaru $(\exists x_1)\dots(\exists x_n) \;p(x_1,\dots,x_n)=q(x_1,\dots,x_n)$
\begin{itemize}
\item kde $p$ a $q$ jsou polynomy s přirozenými koeficienty
\item věta: neexistuje algoritmus, který by pro danou dvojici polynomů s přirozenými koeficienty rozhodl, zda mají přirozené řešení, tj. zda platí $\underline{\mathbb N}\vDash\varphi$
\end{itemize}
\item věta: neexistuje algoritmus, který by pro danou vstupní formuli $\varphi$ rozhodl, zda je logicky platná (tedy zda je tautologie)
\end{itemize}
\item příklady rozhodnutelných teorií
\begin{itemize}
\item teorie čisté rovnosti (bez axiomů, v jazyce $L=\braket{}$ s rovností)
\item teorie unárního predikátu (bez axiomů, v jazyce $L=\braket{U}$ s rovností, kde $U$ je unární relační symbol)
\item teorie hustých lineárních uspořádání DeLO*
\item teorie komutativních grup
\item teorie Booleových algeber
\item teorie následníka s nulou
\item Presburgerova aritmetika
\end{itemize}
\item příklady nerozhodnutelných teorií
\begin{itemize}
\item Robinsonova aritmetika (a také Peanova aritmetika a ZFC, což jsou její extenze)
\item teorie (obecných) grup
\end{itemize}
\end{itemize}
\end{itemize}
